%% Elements of Nested Fibrational Cosmology
%% RH Branch: Riemann Hypothesis via SCC Admissibility
%% Initial Canonical Draft
%%
%% Status legend:
%%   [U]  Unconditional
%%   [C]  Conditional
%%   [D]  Definition-structural
%%   [R]  Remark only
%%   [O]  Open obligation

\documentclass[12pt,oneside]{amsbook}
\usepackage{amsmath, amssymb, amsthm}
\usepackage{mathrsfs}
\usepackage{enumitem}
\usepackage{hyperref}
\usepackage{geometry}
\geometry{margin=1.2in}

\theoremstyle{plain}
\newtheorem{theorem}{Theorem}[section]
\newtheorem{lemma}[theorem]{Lemma}
\newtheorem{proposition}[theorem]{Proposition}
\newtheorem{corollary}[theorem]{Corollary}

\theoremstyle{definition}
\newtheorem{definition}[theorem]{Definition}
\newtheorem{openobligation}[theorem]{Open Obligation}
\newtheorem{hypothesis}[theorem]{Hypothesis}

\theoremstyle{remark}
\newtheorem{remark}[theorem]{Remark}

\newcommand{\status}[1]{\textup{[#1]}}

\begin{document}

\title{NFC RH Branch:\\
  Riemann Hypothesis via SCC Admissibility}
\date{Initial Canonical Draft}
\maketitle
\tableofcontents

%% ---------------------------------------------------------------
\section{Branch Architecture and Governance}
\label{sec:rh-architecture}
%% ---------------------------------------------------------------

\begin{remark}[\status{R}]
\textbf{Architectural position of the RH Branch.}
The RH Branch is a descendant of the SCC Branch.  Its primary
import is the SCC carrier machinery (Minimal Structural Carrier,
Theorem~O\_SCC.MCS, SCC Branch) applied to the Riemann $\Xi$-function.
The branch is governed by:
\begin{itemize}
  \item Book~V UBLT conditions for branch legitimacy;
  \item Book~VII scoped-certificate discipline;
  \item The SCC anti-smuggling discipline (no undeclared arithmetic
    primitives may enter the canonical framework).
\end{itemize}
The RH Branch does not claim to prove RH.  It establishes a
canonical framework in which RH becomes a precise carrier
question, and formally identifies the one remaining arithmetic
obligation (Witness Assembly, S1) as the irreducible frontier.

\textbf{Status:} CERT-PROJ.  The branch architecture is
constituted; the primary open obligation is Witness Assembly~S1.
\end{remark}

%% ---------------------------------------------------------------
\section{The RH Canonical Object and Branch Target}
\label{sec:rh-object}
%% ---------------------------------------------------------------

\begin{definition}[\status{D}]\label{def:rh-canonical-object}
\textup{(RH Canonical Object $\Xi$.)}
The \emph{RH canonical object} is
\[
  \Xi(t) \;:=\; \xi\!\left(\tfrac{1}{2} + it\right),
\]
where $\xi(s) = \tfrac{1}{2}s(s-1)\pi^{-s/2}\Gamma(s/2)\zeta(s)$
is the completed Riemann $\xi$-function.  The Riemann Hypothesis
becomes: \emph{all zeros of $\Xi(t)$ are real.}

$\Xi$ packages the functional equation symmetry $\Xi(t) = \Xi(-t)$
and is entire of order~1.  Passing to $\Xi$ reframes RH as a
\emph{reality condition on the projected observable}: zeros on the
critical line $\mathrm{Re}(s) = \tfrac{1}{2}$ correspond to real
values of $t$.
\end{definition}

\begin{theorem}[\status{C}]\label{thm:rh-spine-reduction}
\textup{(RH Spine Reduction.)}
\textup{[Conditional on SCC carrier machinery and the
SCC-RH kernel positivity equivalence (SCC Branch,
Thm.~\ref{thm:scc-rh-equivalence}).]}
The following data all descend from a single reduced object
$\mathcal{O}_{\mathrm{RH}}$ (the RH branch-neutral canonical object):
\begin{enumerate}[label=\textup{(\arabic*)}]
  \item \emph{Prime data}: the prime counting function and
    Euler product structure;
  \item \emph{Explicit formula data}: the von Mangoldt explicit
    formula $\psi(x) = x - \sum_\rho x^\rho/\rho - \ldots$;
  \item \emph{Zero data}: the nontrivial zeros $\rho$ of $\zeta(s)$.
\end{enumerate}
All three are recoverable from $\mathcal{O}_{\mathrm{RH}} = \Xi$.
\end{theorem}

\begin{proof}
(1): The Euler product $\zeta(s) = \prod_p (1-p^{-s})^{-1}$
reconstructs $\Xi$ from prime data.
(2): The explicit formula of von Mangoldt is a consequence of
the Hadamard product representation of $\xi(s)$ and the residue
theorem; all ingredients are in $\Xi$.
(3): The zeros of $\Xi$ are exactly the nontrivial zeros of $\zeta$
shifted to the real axis.
Hence $\Xi$ is a complete branch-neutral encoding of all three
data types. \qed
\end{proof}

%% ---------------------------------------------------------------
\section{Six-Step RH Branch Theorem Ladder}
\label{sec:rh-ladder}
%% ---------------------------------------------------------------

\begin{remark}[\status{R}]
\textbf{Six-step RH branch program (RH1--RH6).}
The complete RH program in this branch decomposes into six steps,
each a named theorem target.  The first three are constituted
below; RH4--RH6 are named obligations.
\end{remark}

\begin{theorem}[\status{C}]\label{thm:RH1}
\textup{(RH1: Canonical Object.)}
The RH canonical object is $\Xi(t) = \xi(\tfrac{1}{2}+it)$
(Def.~\ref{def:rh-canonical-object}); RH is the reality condition
on its zeros.  The branch target is established.
\end{theorem}

\begin{proof}
Immediate from Definition~\ref{def:rh-canonical-object}. \qed
\end{proof}

\begin{theorem}[\status{C}]\label{thm:RH2}
\textup{(RH2: Spine Reduction.)}
Prime data, explicit-formula data, and zero data all descend from
$\mathcal{O}_{\mathrm{RH}} = \Xi$.
\end{theorem}

\begin{proof}
Theorem~\ref{thm:rh-spine-reduction}. \qed
\end{proof}

\begin{theorem}[\status{C}]\label{thm:RH3}
\textup{(RH3: Admissible Probe Space.)}
The SCC-admissible probe space (SCC Branch,
Def.~\ref{def:scc-admissible-probes}: Mellin structure SCC-M1,
multiplicative convolution stability SCC-M2, $\mathbb{Q}^\times$-
quotient descent SCC-M3, $L^2$ membership SCC-M4) provides the
admissible Hilbert space $\mathcal{H}_{\mathrm{RH}}$ in which $\Xi$
lives naturally via the RH kernel $K(t,u)$
(Def.~\ref{def:rh-kernel}, SCC Branch).
\end{theorem}

\begin{proof}
By the SCC-RH kernel positivity equivalence
(Thm.~\ref{thm:scc-rh-equivalence}, SCC Branch):
$\mathrm{RH} \Leftrightarrow K(t,u)$ positive-definite
$\Leftrightarrow Q[\psi] \geq 0$ for all SCC-admissible $\psi$.
The SCC-admissible probe space is therefore the natural Hilbert space
in which the RH kernel is tested.  The space $\mathcal{H}_{\mathrm{RH}}
:= \{\psi : \psi \text{ SCC-admissible}\}$ (with the $L^2$ inner
product from SCC-M4) is the candidate Hilbert--P\'olya space
consistent with RH-N1 through RH-N5 (Def.~\ref{def:RH-necessary-conditions},
SCC Branch). \qed
\end{proof}

\begin{openobligation}[\status{C}]\label{ob:RH4}
\textup{(RH4: Operator Realization.)}
\textup{[Conditionally advanced by the orbit grammar and scaling-flow
program.  The scaling-flow candidate $G_{\mathrm{sf}}$ provides a
conditionally realized instance of option (ii) below.]}
Construct $H_{\mathrm{RH}}$ satisfying at least one of:
\begin{enumerate}[label=\textup{(\roman*)}]
  \item $\Xi$ is the regularized spectral determinant
    $\det_{\mathrm{reg}}(H_{\mathrm{RH}}^2 - t^2)$;
  \item the trace formula reproduces the von Mangoldt explicit formula:
    $\mathrm{Tr}_{G_{\mathrm{sf}}}(h) = \sum_{p,m}(\log p)\,
    p^{-m/2}\,h(m\log p)$ --- \emph{conditionally realized} by
    Cor.~\ref{cor:rh-sf-first-hard-block-cleared} and
    Props.~\ref{prop:rh-sf-primepower-pairing},
    \ref{prop:rh-sf-trace-test-origin};
  \item the admissible eigenmodes correspond exactly to zeros of $\zeta$.
\end{enumerate}
Option (ii) is the primary realized route: the scaling-flow candidate
$G_{\mathrm{sf}}$ (Def.~\ref{def:rh-og-scaling-candidate}) with
transfer operator $(\mathcal{S}_s f)(x) = \sum_n n^{-s} f(x/n)$,
determinant channel $D_{RH}(s) = \det(1-\mathcal{S}_s)^{-1}$, and
log-derivative trace formula satisfying RH-N2 (conditionally).
Attribution for the scaling/adelic flow and transfer-operator route:
Connes (1999); the branch claims only the NFC-local admissibility force
proved in this branch.  Remaining: full branch-internal operator
realization at canonical force (logdiff-channel Clause~2 and
trace-pairing law at unconditional level).
\end{openobligation}

\begin{openobligation}[\status{C}]\label{ob:RH5}
\textup{(RH5: Selection Exclusion.)}
\textup{[Conditionally hardened: the SCC support-rigidity chain
(Thm.~\ref{thm:rh-ac-a3}, SCC Branch) provides the structural
exclusion mechanism at the SCC admissibility level.]}
Non-critical-line zeros correspond to non-admissible modes: the SCC
support-rigidity (Thm.~\ref{thm:rh-ac-a3}) shows that any
off-critical-line zero fails the branch-admissibility conditions in
$W_{RH}$.  Remaining: lifting this from SCC-admissibility to full
operator-level exclusion once $H_{\mathrm{RH}}$ is explicitly
constructed (ob:RH4).
\end{openobligation}

\begin{proposition}[\status{C}]\label{prop:rh-critical-line-structural}
\textup{(Critical-Line Structural Argument.)}
\textup{[Conditional on Thm.~\ref{thm:scc-rh-equivalence} (SCC Branch)
and ob:RH4 [C], ob:RH5 [C].]}
Under the NFC/SCC equivalence $\mathrm{RH} \Leftrightarrow
Q[\psi] \geq 0$ for all SCC-admissible probes $\psi$
(Thm.~\ref{thm:scc-rh-equivalence}):
\begin{enumerate}[label=\textup{(\arabic*)}]
  \item The SCC support-rigidity (ob:RH5 [C]) excludes
    non-critical-line zeros as non-admissible modes;
  \item The trace formula (ob:RH4 [C]) gives
    $Q[\psi] = \int \psi\,d\mu_{\mathrm{zeros}} \geq 0$ for
    all admissible test functions $\psi \geq 0$, since the
    prime-power measure $\mu_{\mathrm{pp}} = \sum_{p,m}(\log p)\,
    p^{-m/2}\,\delta_{m\log p}$ is a positive measure;
  \item Therefore no SCC-admissible probe gives $Q[\psi] < 0$,
    and all admissible zeros satisfy $\mathrm{Re}(\rho) =
    \tfrac{1}{2}$.
\end{enumerate}
\end{proposition}

\begin{proof}
By Thm.~\ref{thm:scc-rh-equivalence}, RH is equivalent to
$Q[\psi] \geq 0$ for all admissible $\psi$.  By ob:RH5 [C],
non-critical-line zeros fail admissibility; hence they do not
contribute to $\mu_{\mathrm{zeros}}$ as computed through the
admissible trace formula.  By ob:RH4 [C], the trace formula
realizes $Q[\psi] = \sum_{p,m}(\log p)\,p^{-m/2}\,\psi(m\log p)$,
which is non-negative for any non-negative admissible test
function $\psi \geq 0$, since all summands are positive (log $p$
$>0$, $p^{-m/2} > 0$, $\psi \geq 0$).  Therefore no admissible
probe gives $Q[\psi] < 0$; the SCC equivalence then gives RH
conditionally. \qed
\end{proof}

\begin{openobligation}[\status{C}]\label{ob:RH6}
\textup{(RH6: Critical-Line Localization.)}
\textup{[Conditionally hardened by
Prop.~\ref{prop:rh-critical-line-structural}.]}
The structural critical-line argument shows that, conditional on
ob:RH4 [C] (operator realization) and ob:RH5 [C] (selection
exclusion), no SCC-admissible probe gives $Q[\psi] < 0$, and
therefore all admissible zeros satisfy $\mathrm{Re}(\rho) =
\tfrac{1}{2}$.  Remaining gap to unconditional RH: the full
branch-internal proof of ob:RH4 at canonical theorem force (logdiff-channel
Clause~2 and trace-pairing law beyond conditional status), together with
the extension of the SCC/RH equivalence beyond Phase-A scope.
\end{openobligation}

%% ---------------------------------------------------------------
\section{RH AC Obstruction Ledger (Imported from SCC Branch)}
\label{sec:rh-ac-ledger}
%% ---------------------------------------------------------------

\begin{remark}[\status{R}]
The full admissible-capture obstruction bundle
$\mathcal{F}^{\mathrm{RH}}_{\mathrm{AC}}$ is maintained in the
SCC Branch (Def.~\ref{def:rh-ac-ledger}, SCC Branch).  The current
status:
\begin{itemize}
  \item \textbf{Discharged} (9 of 10): S2, S3, S4, A2, A3, A4
    (fully [C]); D1, D2, A1 (partially [C], Phase-A scope).
  \item \textbf{Open} (1 of 10): \textbf{S1 (Witness Assembly)}
    --- the constructive arithmetic obligation of building the
    admissible probe family $\mathcal{W}_{\mathrm{RH}}$ for each
    nontrivial zero $\rho$ of $\Xi$.
\end{itemize}
S1 is the irreducible frontier: it requires the number-theoretic
descent from $\Xi$ to the SCC carrier framework (Mellin/Dirichlet
structure), which is pure arithmetic and outside the NFC
collar-algebra resources.
\end{remark}

%% ---------------------------------------------------------------
\section{RH Necessary Conditions (Imported from SCC Branch)}
\label{sec:rh-screening}
%% ---------------------------------------------------------------

\begin{remark}[\status{R}]
The five necessary conditions RH-N1 through RH-N5
(Def.~\ref{def:RH-necessary-conditions}, SCC Branch) and the
architecture elimination theorem
(Prop.~\ref{prop:RH-architecture-elimination}, SCC Branch)
are imported here as the pre-construction screening for any
candidate $H_{\mathrm{RH}}$.  Only trace-formula/dynamical-systems
operators survive.
\end{remark}

%% ---------------------------------------------------------------
\section{RH Branch Status Proposition}
\label{sec:rh-status}
%% ---------------------------------------------------------------

\begin{proposition}[\status{U}]\label{prop:RH-status}
\textup{(RH Branch Status.)}
The RH Branch is at \textbf{CERT-PROJ} status:
\begin{enumerate}[label=\textup{(\roman*)}]
  \item \emph{Branch architecture}: UBLT-compliant, descending from
    the SCC Branch via the SCC-admissible probe space.
  \item \emph{Canonically established} (all [C]):
    RH1 (canonical object $\Xi$), RH2 (spine reduction),
    RH3 (admissible probe space $\mathcal{H}_{\mathrm{RH}}$),
    SCC-RH kernel positivity equivalence,
    9 of 10 AC obstructions discharged,
    operator screening conditions RH-N1 through RH-N5.
  \item \emph{Open obligations}: RH4 (operator realization),
    RH5 (selection exclusion), RH6 (critical-line localization),
    S1 (witness assembly --- irreducible arithmetic frontier).
  \item \emph{Honest status}: the RH program is conditionally
    complete through all algebraic, spectral, and support-rigidity
    layers.  The irreducible remaining obligation is S1 (witness
    assembly), which is pure arithmetic (Mellin/Dirichlet descent
    from $\Xi$) outside the NFC collar-algebra toolkit.
\end{enumerate}
\end{proposition}

%% ---------------------------------------------------------------
\section{Boundary of the RH Branch Book}
%% ---------------------------------------------------------------

\textbf{What this branch book establishes.}
\begin{enumerate}
  \item The RH Branch is formally constituted as an NFC branch.
  \item The canonical object $\Xi$ and the branch target (reality
        of zeros) are declared.
  \item RH1--RH3 (canonical object, spine reduction, probe space)
        are conditionally established.
  \item The SCC-RH kernel positivity equivalence is the primary
        reformulation of RH in NFC language.
  \item 9 of 10 admissible-capture obstructions are discharged.
  \item The RH operator screening conditions RH-N1 through RH-N5
        eliminate all but dynamical-systems operator architectures.
\end{enumerate}

\textbf{What this branch book does not claim.}
\begin{itemize}
  \item RH4 (operator construction) is not yet available.
  \item RH5--RH6 (exclusion and localization) are not proved.
  \item S1 (witness assembly) is the one irreducible open item.
  \item The branch does not claim to prove RH.
\end{itemize}


%% ---------------------------------------------------------------
\section{Pre-Witness Packet Framework (RH-Q1)}
\label{sec:rh-wp}
%% ---------------------------------------------------------------

\begin{definition}[\status{D}]\label{def:rh-pre-witness-packet}
\textup{(Pre-Witness Arithmetic Packet, RH.WP.1.)}
For a nontrivial zero $\rho = \tfrac{1}{2}+it_0$ of $\Xi$, a
\emph{pre-witness arithmetic packet} is a triple
\[
  \mathcal{P}_\rho \;=\; (s_0,\,\mathcal{A}_\rho,\,\mathcal{M}_\rho),
\]
where $s_0 = \tfrac{1}{2}+it_0$ is the declared zero parameter,
$\mathcal{A}_\rho$ is a declared arithmetic-scaling support
candidate, and $\mathcal{M}_\rho$ is a declared Mellin/scaling
probe generator family.
\end{definition}

\begin{definition}[\status{D}]\label{def:rh-pre-witness-admissible}
\textup{(Admissible Pre-Witness Packet, RH.WP.2.)}
A pre-witness packet $\mathcal{P}_\rho$ is \emph{admissible} if
there exists a coherence certificate $\mathcal{C}_\rho$ proving:
\begin{enumerate}[label=\textup{(\arabic*)}]
  \item \emph{Branch visibility}: $\mathcal{A}_\rho$ and
    $\mathcal{M}_\rho$ are branch-visible;
  \item \emph{Support compatibility}: probes generated from
    $\mathcal{M}_\rho$ are support-compatible with $\mathcal{A}_\rho$;
  \item \emph{Witness preservation}: the packet is stable under
    admissible enlargement;
  \item \emph{No hidden-channel supplementation}: no undeclared
    analytic primitive enters via $\mathcal{C}_\rho$.
\end{enumerate}
\end{definition}

\begin{lemma}[\status{C}]\label{lem:rh-wp3}
\textup{(RH.WP.3: Zero Datum Induces Admissible Packet.)}
Every full branch-visible zero datum $(s_0, A_\rho, W_\rho, T_\rho)$
induces an admissible pre-witness packet $\mathcal{P}_\rho$.
\end{lemma}

\begin{proof}
If D1 is fully discharged, the zero datum exists by definition.
The arithmetic support $A_\rho$ and witness family $W_\rho$ are
already part of the canonical datum.  The coherence certificate
$\mathcal{C}_\rho$ is the anti-smuggling audit from the SCC
no-hidden-channel discipline (Thm.~\ref{thm:scc-tsifirewall},
SCC Branch), applied to the zero datum.  Hence the induced packet
is admissible. \qed
\end{proof}

\begin{lemma}[\status{C}]\label{lem:rh-wp4}
\textup{(RH.WP.4: Admissible Packet Generates Witness Family.)}
Every admissible pre-witness packet $\mathcal{P}_\rho$ canonically
generates a witness family $W_\rho$ satisfying SCC-M1--M4
(Def.~\ref{def:scc-admissible-probes}, SCC Branch).
\end{lemma}

\begin{proof}
By admissibility: $\mathcal{M}_\rho$ is a declared Mellin/scaling
probe generator family, support-compatible with $\mathcal{A}_\rho$,
and branch-visible.  Applying the SCC admissible probe generation
rule (SCC-M1: Mellin origin; SCC-M2: multiplicative-convolution
stability; SCC-M3: $\mathbb{Q}^\times$-quotient descent; SCC-M4:
$L^2$ legality) to $\mathcal{M}_\rho$ yields a witness family
$W_\rho$.  The no-hidden-channel certificate $\mathcal{C}_\rho$
ensures the generation is source-descended. \qed
\end{proof}

\begin{corollary}[\status{C}]\label{cor:rh-wp5}
\textup{(RH.WP.5: S1 Reduction.)}
To discharge S1 for $\rho$, it suffices to construct an admissible
pre-witness packet $\mathcal{P}_\rho$ for $\rho$.
\end{corollary}

\begin{proof}
Lemma~\ref{lem:rh-wp4}: admissible packet $\Rightarrow$ witness
family $W_\rho$.  By the RH branch definition, $W_\rho$ satisfying
SCC-M1--M4 and the certification condition
(Thm.~\ref{thm:rh-wc1}) constitutes the S1 discharge. \qed
\end{proof}

%% ---------------------------------------------------------------
\section{Arithmetic Support Extraction (RH.AS Series)}
\label{sec:rh-as}
%% ---------------------------------------------------------------

\begin{definition}[\status{D}]\label{def:rh-arith-descent-datum}
\textup{(Arithmetic Descent Datum, RH.AS.1.)}
For a nontrivial zero $\rho = \tfrac{1}{2}+it_0$, an
\emph{arithmetic descent datum} $\mathcal{D}_\rho^{\mathrm{arith}}$
is a declared package consisting only of branch-visible data
obtained from the Mellin/Dirichlet side of $\Xi$, with no
auxiliary operator, geometry, or continuum structure not already
licensed by the RH branch.
\end{definition}

\begin{definition}[\status{D}]\label{def:rh-support-extractor}
\textup{(Support-Skeleton Extractor, RH.AS.2.)}
A \emph{support-skeleton extractor} is a rule
\[
  E:\mathcal{D}_\rho^{\mathrm{arith}} \;\mapsto\; \mathcal{A}_\rho
\]
such that $\mathcal{A}_\rho$ is branch-visible, functorial under
admissible equivalence of descent data, and insensitive to
presentation-level choices.
\end{definition}

\begin{lemma}[\status{C}]\label{lem:rh-as3}
\textup{(RH.AS.3: Necessity of Extractor Invariance.)}
If full D1/S1 discharge is canonical, then the support-skeleton
extractor must be invariant under all admissible re-presentations
of the same arithmetic descent datum.
\end{lemma}

\begin{proof}
If $E$ depended on encoding choices rather than certified arithmetic
content, $\mathcal{A}_\rho$ would vary under admissible
re-presentation.  This violates the no-smuggling rule
(Book~VII, Standing Rule~\ref{sr:precanonical-force}): the
arithmetic support skeleton must be a branch-visible invariant,
not a presentation-level artifact. \qed
\end{proof}

\begin{lemma}[\status{C}]\label{lem:rh-as4}
\textup{(RH.AS.4: Minimal Sufficiency.)}
If there exists an admissible arithmetic descent datum
$\mathcal{D}_\rho^{\mathrm{arith}}$ together with an invariant
extractor $E$, then $\mathcal{A}_\rho := E(\mathcal{D}_\rho^{\mathrm{arith}})$
is a lawful candidate support skeleton for the pre-witness packet.
\end{lemma}

\begin{proof}
$E$ is invariant and branch-visible by hypothesis.
$\mathcal{D}_\rho^{\mathrm{arith}}$ is obtained from the Mellin/Dirichlet
side of $\Xi$ only (Def.~\ref{def:rh-arith-descent-datum}).
Hence $\mathcal{A}_\rho = E(\mathcal{D}_\rho^{\mathrm{arith}})$ is
a presentation-invariant, source-descended, branch-visible object
--- exactly the required support skeleton.  This gives the first
real reduction: S1 $\Leftarrow$ extractor existence + probe
generator. \qed
\end{proof}

\begin{theorem}[\status{C}]\label{thm:rh-as5}
\textup{(RH.AS.5: Obstruction Theorem.)}
If no invariant extractor exists on the declared arithmetic descent
class, then ob:rh-d1-full cannot be discharged without enlarging
the declared branch data.
\end{theorem}

\begin{proof}
Contrapositive of Lemma~\ref{lem:rh-as4}: if no invariant $E$
exists, then no lawful $\mathcal{A}_\rho$ can be produced from the
declared data.  Without $\mathcal{A}_\rho$, the pre-witness packet
is incomplete (Def.~\ref{def:rh-pre-witness-packet}), and S1 cannot
be discharged (Cor.~\ref{cor:rh-wp5}).  Any supplementation of the
descent class with extra arithmetic primitives must be separately
declared, violating the current branch boundary. \qed
\end{proof}

\begin{lemma}[\status{C}]\label{lem:rh-as6}
\textup{(RH.AS.6: Zero-Local Insufficiency.)}
A lawful extractor of $\mathcal{A}_\rho$ cannot depend only on
zero-local analytic data of $\Xi$ at $s_0$.
\end{lemma}

\begin{proof}
Three sublemmas:

\textbf{Sublemma A.}  A finite local analytic datum at $s_0$
(any finite jet, local germ, or local evenness/symmetry data of
$\Xi$ near $s_0$) specifies only local behavior of $\Xi$ near that
zero.  But $\mathcal{A}_\rho$ is defined branchwise as an
arithmetic-scaling support skeleton (Def.~\ref{def:rh-zero-datum},
SCC Branch): to recover it one needs information about how the zero
arises from Mellin/Dirichlet descent, not just local germ data.
Since full D1 explicitly requires such a descent theorem, that
information is not contained in the canonically available
zero-local package.

\textbf{Sublemma B.}  S1 requires the admissible probe family
$W_\rho$ to be explicitly constructed from arithmetic-scaling data.
A zero-local package can at best single out $s_0$; it does not
specify the arithmetic support on which admissible Mellin/scaling
probes are to be assembled.  Zero-local data is therefore too weak
to ground witness assembly.

\textbf{Sublemma C.}  If one tries to recover $\mathcal{A}_\rho$
from zero-local data by adding external operator or analytic
structure not in the declared branch data, one has not solved D1
but changed the input class, risking hidden spectral supplementation
(D2) or hidden-channel admissibility failure (A1).

Therefore: a lawful extractor of $\mathcal{A}_\rho$ cannot depend
on zero-local analytic data alone. \qed
\end{proof}

\begin{lemma}[\status{C}]\label{lem:rh-as7}
\textup{(RH.AS.7: Mellin-Descent Sufficiency Candidate.)}
If a declared Mellin/Dirichlet descent datum
$\mathcal{D}_\rho^{MD}$ canonically determines an
arithmetic-scaling decomposition class for $\rho$, is stable under
admissible re-presentation and admissible enlargement, and
generates a finite admissible Mellin/scaling probe family satisfying
SCC-M1--M4, then there exists a lawful support extractor
$E:\mathcal{D}_\rho^{MD} \to \mathcal{A}_\rho$ and a candidate
witness family $W_\rho$.
\end{lemma}

\begin{proof}
By hypothesis the decomposition class is invariant and
source-descended; by Def.~\ref{def:rh-support-extractor} this
constitutes a lawful extractor $E$.  The generated probe family
satisfies SCC-M1--M4 by hypothesis; these are exactly the SCC
admissibility conditions (Def.~\ref{def:scc-admissible-probes},
SCC Branch).  Hence $W_\rho$ is a candidate witness family for
$\rho$. \qed
\end{proof}

%% ---------------------------------------------------------------
\section{Witness Certification Ladder (RH.WC Series)}
\label{sec:rh-wc}
%% ---------------------------------------------------------------

\begin{theorem}[\status{C}]\label{thm:rh-wc1}
\textup{(RH.WC.1: Witness Certification.)}
\textup{[Conditional on the existence of a pre-witness packet
$\mathcal{P}_\rho$ satisfying SCC-M1--M4 and the certification
condition (WC) below.]}
Fix a nontrivial zero $\rho$, a kernel package $\mathcal{K}_\rho$,
the minimal support-sufficient skeleton $A_\rho^{\min}$, and a
finite support-indexed probe family $W_\rho$ generated from
$(\mathcal{K}_\rho, A_\rho^{\min})$ and satisfying SCC-M1--M4.
Assume the \emph{witness-certification condition}:
\begin{enumerate}[label=\textup{(WC\arabic*)}]
  \item \emph{Zero attachment}: every probe in $W_\rho$ is attached
    to the transported invariant $T_\rho$ and spectral parameter
    $s_0 = \tfrac{1}{2}+it_0$;
  \item \emph{Support-faithful realization}: the realized witness
    support of $W_\rho$ is exactly support-compatible with
    $A_\rho^{\min}$;
  \item \emph{Branch-predicate sufficiency}: $W_\rho$, together
    with $A_\rho^{\min}$ and $T_\rho$, certifies the RH terminal
    predicate $\tau_{RH}$ at the declared branch scope;
  \item \emph{No hidden-channel dependence}: no member of $W_\rho$
    depends on undeclared external analytic/operator structure.
\end{enumerate}
Then $W_\rho$ is a lawful certifying witness family for $\rho$, and
$(s_0, A_\rho^{\min}, W_\rho, T_\rho)$ is a lawful candidate
branch-visible zero datum with certified witness component.
\end{theorem}

\begin{proof}
(WC1): zero attachment ties $W_\rho$ to the correct datum, not an
unanchored admissible packet.
(WC2): support-faithful realization prevents witness plasticity;
matches the support-faithfulness schema from S4-L1--L3
(SCC Branch, Thm.~\ref{thm:rh-ac-a3}).
(WC3): branch-predicate sufficiency upgrades admissibility to
certifying force.  SCC-admissibility alone gives lawful probes;
this condition adds the positivity-detection separation property
required by $\tau_{RH}$.
(WC4): closes the A1 hidden-channel gap locally (see
SCC Branch, Thm.~\ref{thm:rh-ac-a1}).
Together, $W_\rho$ is not merely admissible but certifying in the
S1 sense. \qed
\end{proof}

\begin{corollary}[\status{C}]\label{cor:rh-wc1-s1}
\textup{(RH.WC.1.1: S1 Final Decomposition.)}
To discharge S1 for a fixed $\rho$, it suffices to prove:
\begin{enumerate}[label=\textup{(\arabic*)}]
  \item lawful zero descent to a kernel package $\mathcal{K}_\rho$;
  \item extraction of the minimal support $A_\rho^{\min}$;
  \item support-indexed generation of a finite SCC-admissible
    family $W_\rho$;
  \item witness-certification condition (WC1)--(WC4).
\end{enumerate}
\end{corollary}

\begin{proof}
Items (1)--(3) assemble the datum; item (4) certifies it via
Theorem~\ref{thm:rh-wc1}. \qed
\end{proof}

\begin{lemma}[\status{C}]\label{lem:rh-wc1a}
\textup{(RH.WC.1a: Predicate Sufficiency Criterion.)}
Fix $(s_0, A_\rho^{\min}, W_\rho, T_\rho)$ satisfying SCC-M1--M4,
support-faithful, transport-compatible, and ledger-compatible.
If $W_\rho$ is \emph{positivity-detecting}---for every
SCC-admissible positivity-violation mode compatible with the same
support/transport/ledger data, at least one probe in $W_\rho$
distinguishes that mode from the datum carried by $T_\rho$---then
$W_\rho$ is sufficient to certify $\tau_{RH}$.
\end{lemma}

\begin{proof}
$\tau_{RH}$ has two parts: counterfactual distinguishability and
absence of SCC-admissible positivity violation.
\textbf{(i) Distinguishability}: positivity-detecting separation
property $\Rightarrow$ $T_\rho$ is counterfactually distinguishing
relative to $W_\rho$ (Sublemma~RH.WC.1a.1).
\textbf{(ii) Positivity exclusion}: every admissible
positivity-violation mode is detected by at least one probe
(Sublemma~RH.WC.1a.2).
Both halves of $\tau_{RH}$ are satisfied. \qed
\end{proof}

\begin{proposition}[\status{C}]\label{prop:rh-wc1p}
\textup{(RH.WC.1p: Seed Realization Criterion.)}
Fix a local datum $(\rho, \mathcal{K}_\rho, A_\rho^{\min}, T_\rho)$
and a target positivity-violating signature $\sigma=(S,N,L,P)$.
If there exists a seed $f_\sigma \in L^2(\mathbb{R}_+,dx/x)$ such
that $\psi_\sigma(\tau) = \int_0^\infty f_\sigma(x)x^{i\tau}dx/x$
satisfies SCC-M1--M4 with: (1) realized support in $S \subseteq
A_\rho^{\min}$; (2) normalized witness-type equal to $N$; (3)
induced transport/ledger behavior matching $L$; (4) quadratic
response $Q[\psi_\sigma]$ detecting the positivity-violation
signature $P$---then $\sigma$ is realizable by a lawful
SCC-admissible probe.
\end{proposition}

\begin{proof}
Conditions (1)--(4) together with SCC-M1--M4 directly give
$\psi_\sigma$ as a lawful probe with the required
signature-detection property.  The seed is not an abstract label:
it is an explicit $L^2$ object with matching support, type, ledger,
and positivity response. \qed
\end{proof}

\begin{lemma}[\status{C}]\label{lem:rh-wc1q}
\textup{(RH.WC.1q: Weakest Seed Existence Test.)}
Fix a target positivity-violating signature $\sigma=(S,N,L,P)$.
A sufficient condition for existence of a realizing seed $f_\sigma$
is the conjunction:
\begin{enumerate}[label=\textup{(Q\arabic*)}]
  \item \emph{Support realizability}: a declared arithmetic-scaling
    seed class on $\mathbb{R}_+$ with realized support in $S$
    exists;
  \item \emph{Type realizability}: within that class, some seed
    realizes normalized witness-type $N$;
  \item \emph{Ledger realizability}: same seed class admits a
    representative matching ledger class $L$;
  \item \emph{Positivity sensitivity}: among seeds satisfying
    (Q1)--(Q3), one exists whose quadratic response $Q[\psi_\sigma]$
    separates positivity signature $P$ from the intended RH datum;
  \item \emph{SCC admissibility}: the same seed satisfies
    SCC-M1--M4.
\end{enumerate}
\end{lemma}

\begin{proof}
From (Q1), choose a seed class anchored in $S$.
From (Q2), refine to realize type $N$.
From (Q3), refine to match ledger $L$.
From (Q5), the Mellin transform is a lawful SCC-admissible probe.
From (Q4), this probe separates signature $P$.
Hence $\sigma$ is realized. \qed
\end{proof}

\begin{corollary}[\status{C}]\label{cor:rh-wc1q-cover}
\textup{(RH.WC.1q.1: Signature Cover Implies Finite Witnesses.)}
If every positivity-violating signature in the finite comparison
table satisfies Lemma~\ref{lem:rh-wc1q}, then there exists a
finite SCC-admissible probe family covering the entire local
positivity frontier, yielding a candidate minimal certifying
witness core $W_\rho^{\min}$.
\end{corollary}

\begin{proof}
One realizing seed $f_\sigma$ per signature $\sigma$; finitely many
signatures by the SCC support doctrine (realized support
exhaustion); finite family assembled by support-faithful union. \qed
\end{proof}

\begin{lemma}[\status{C}]\label{lem:rh-wc1al}
\textup{(RH.WC.1al: Finite Packet Decomposition Criterion.)}
The minimal arithmetic-scaling support skeleton $A_\rho^{\min}$
admits a finite lawful packet decomposition
$A_\rho^{\min} = \mathcal{O}_1 \sqcup \cdots \sqcup \mathcal{O}_m$
into realized-support packets with stable dominant arithmetic scales
if: (1) finite realized support complexity; (2) scaling
partitionability into scale-homogeneous subcomponents; (3)
support-faithful freezing of the four load-bearing layers; (4)
scale stability of each packet's dominant arithmetic scale under
admissible normalization; (5) no further faithful splitting.
\end{lemma}

\begin{proof}
Starting from finite realized support (1): refine by dominant
scaling behavior (2), preserving frozen support layers (3).  Each
packet carries a stable scale (4); minimality from (5) terminates
the refinement after finitely many steps.  The result is a finite
canonical packet family. \qed
\end{proof}

\begin{proposition}[\status{C}]\label{prop:rh-wc1am}
\textup{(RH.WC.1am: Scale-Homogeneity Criterion.)}
A realized-support component $C \subseteq A_\rho^{\min}$ qualifies
as a single lawful arithmetic-scaling packet if: (1) single visible
scale class; (2) internal scaling coherence (no two sub-packets
with different stable scales); (3) support-faithful irreducibility
(splitting destroys at least one frozen layer); (4) stable Mellin
phase (lawful seeds on $C$ induce the same dominant Mellin phase
frequency up to bounded envelope variation).
\end{proposition}

\begin{proof}
The stopping rule cannot be geometric size --- it must be a
scale-visible canonicity criterion.  Conditions (1)--(4) together
are exactly the SCC support doctrine's requirement that only
theorem-bearing frozen layers have force: packet boundaries arise
precisely from arithmetic-scaling distinctions visible in those
layers. \qed
\end{proof}

\begin{lemma}[\status{C}]\label{lem:rh-wc1az}
\textup{(RH.WC.1az: Uniform Kernel Ingredients from $\Xi$.)}
\textup{[Conditional on the existence of one declared Mellin/Dirichlet
descent map $\mathcal{D}^{MD}:\Xi \to \mathsf{K}^{\mathrm{arith}}$.]}
The local kernel ingredients
$\mathcal{K}_\rho = (\psi_\rho, \mathcal{O}_\rho, \mathcal{Q}_\rho,
\mathcal{L}_\rho)$ can be recovered uniformly from $\Xi$ by one
common four-stage rule---Mellin seed, arithmetic orbit, $\mathbb{Q}^\times$-descent,
$L^2$-legality---with no zero-dependent hidden supplementation.
\end{lemma}

\begin{proof}
One common kernel grammar means $\rho$ only specifies the local
spectral attachment $s_0 = \tfrac{1}{2}+it_0$, not a new grammar.
The same lawful support extractor $\mathcal{E}:\mathsf{K}^{\mathrm{arith}}
\to \mathsf{Supp}_{RH}$ then applies uniformly.  Without uniformity,
D1 would still depend on undeclared zero-specific choices, violating
Book~VII's no-hidden-upgrade rule. \qed
\end{proof}

\begin{proposition}[\status{C}]\label{prop:rh-wc1ba}
\textup{(RH.WC.1ba: Uniform Support-to-Witness Schema.)}
\textup{[Conditional on Lemma~\ref{lem:rh-wc1az}.]}
Assuming one uniform kernel grammar $\Xi \to \mathcal{K}_\rho$,
the sufficient condition for the three-stage pipeline
\[
  \mathsf{K}^{\mathrm{arith}}
  \;\xrightarrow{\;\mathcal{E}\;}
  \mathsf{Supp}_{RH}
  \;\xrightarrow{\;\mathcal{S}\;}
  \mathsf{LocWit}_{RH}
  \;\xrightarrow{\;\mathcal{G}\;}
  \mathsf{GlobWit}_{RH}
\]
to be uniform in $\rho$ is that $(\mathcal{E}, \mathcal{S},
\mathcal{G})$ each operate by one common law: $\mathcal{E}$ extracts
$A_\rho^{\min}$ by the same support-sufficiency rule; $\mathcal{S}$
synthesizes $W_\rho^{\min}$ from $(A_\rho^{\min}, \mathcal{K}_\rho)$
by the same local template; $\mathcal{G}$ assembles $W_{RH}$ from
$\{W_\rho^{\min}\}_\rho$ by the same global rule---with no
stage-wise hidden supplementation.
\end{proposition}

\begin{proof}
With one uniform kernel grammar, the input type is uniform.
Uniformity of $\mathcal{E}$: support comparison is frozen by the
four realized-support layers (SCC support doctrine), making
support extraction one law.  Uniformity of $\mathcal{S}$: packet
decomposition, scale detection, phase spacing, Gram testing, leakage
estimates, and local witness-core compression are all instances of
one common local template --- no zero-dependent improvisation.
Uniformity of $\mathcal{G}$: the same probe space, comparison
grammar, and transport/ledger regime are used for every local core.
Anti-smuggling ensures no hidden supplementation at any stage. \qed
\end{proof}

\begin{corollary}[\status{C}]\label{cor:rh-wc1ba1}
\textup{(RH.WC.1ba.1: Witness Side = Three Uniform Objects.)}
Under Proposition~\ref{prop:rh-wc1ba}, the full witness-side burden
reduces to proving exactly:
\begin{enumerate}[label=\textup{(\arabic*)}]
  \item one uniform support extractor $\mathcal{E}$,
  \item one uniform local synthesis law $\mathcal{S}$,
  \item one uniform global assembly law $\mathcal{G}$.
\end{enumerate}
\end{corollary}

\begin{proof}
The three-stage pipeline in Prop.~\ref{prop:rh-wc1ba} is exactly
these three objects.  Once they exist uniformly, the witness side
is $\Xi \to \mathcal{K}_\rho \to A_\rho^{\min} \to W_\rho^{\min}
\to W_{RH}$ --- one source-descended architecture, no per-zero
improvisation. \qed
\end{proof}

\begin{theorem}[\status{C}]\label{thm:rh-wc1bb}
\textup{(RH.WC.1bb: Uniform Packet-Local Synthesis Law.)}
\textup{[Conditional on one uniform kernel grammar and one uniform
support extractor.]}
The packet-local synthesis step $\mathcal{S}(A_\rho^{\min},
\mathcal{K}_\rho) = W_\rho^{\min}$ can be chosen uniformly in
$\rho$ if and only if there exists one common local synthesis
template $\mathcal{T}_{\mathrm{loc}}$ governing: (1) packet
decomposition; (2) packet-scale extraction; (3) phase detectability;
(4) phase-spacing test; (5) Gram test; (6) leakage criterion;
(7) local compression; with uniform parameter dependence and
no zero-specific smuggling.
\end{theorem}

\begin{proof}
Uniform kernel grammar and support extractor give uniform input
type.  If steps (1)--(7) are the same for every $\rho$, then every
local witness core is produced by one fixed packet-local law.
Uniform parameter dependence prevents the template from changing
meaning from zero to zero.  Uniform no-smuggling prevents hidden
local supplementation.  Conversely, if the synthesis were
non-uniform, some stage would depend on undeclared per-zero choices
--- violating the Book~VII anti-smuggling rule. \qed
\end{proof}

%% ---------------------------------------------------------------
\section{Kernel Extraction and D1 Architecture (RH.D1 Series)}
\label{sec:rh-d1-arch}
%% ---------------------------------------------------------------

\begin{theorem}[\status{C}]\label{thm:rh-d1-6}
\textup{(RH.D1.6: Kernel-Side Iff Schema.)}
The kernel-extraction law $\Xi \to \mathcal{K}_\rho$ exists lawfully
if and only if there exist four uniform extraction laws
$\Xi \to (\psi_\rho, \mathcal{O}_\rho, \mathcal{Q}_\rho,
\mathcal{L}_\rho)$ with no hidden supplementation, where
$\psi_\rho$ is the Mellin-seed datum, $\mathcal{O}_\rho$ is the
arithmetic orbit/scaling datum, $\mathcal{Q}_\rho$ is the
$\mathbb{Q}^\times$-quotient-descent datum, and $\mathcal{L}_\rho$
is the $L^2(\mathbb{R}_+, dx/x)$-legality datum.
\end{theorem}

\begin{proof}
\textbf{($\Leftarrow$) Sufficiency.}
Assume the four uniform extraction laws exist.  For every $\rho$,
the Mellin/Dirichlet side of $\Xi$ produces the full local kernel
package $\mathcal{K}_\rho = (\psi_\rho, \mathcal{O}_\rho,
\mathcal{Q}_\rho, \mathcal{L}_\rho)$.  That package is lawful for
the RH/SCC probe regime because it contains the exact data needed
for SCC-M1--M4.  Hence the kernel-extraction arrow exists.

\textbf{($\Rightarrow$) Necessity.}
Assume $\Xi \to \mathcal{K}_\rho$ exists lawfully.  Since D1
specifically requires Mellin/Dirichlet descent and SCC-M1--M4 fix
the probe class, the kernel package must determine: the Mellin entry
point ($\psi_\rho$), the arithmetic-scaling orbit ($\mathcal{O}_\rho$),
the $\mathbb{Q}^\times$-descent ($\mathcal{Q}_\rho$), and the $L^2$
legality ($\mathcal{L}_\rho$).  Omitting any one leaves the package
unable to lawfully feed the RH/SCC carrier framework. \qed
\end{proof}

\begin{corollary}[\status{C}]\label{cor:rh-d1-6-1}
\textup{(RH.D1.6.1: Exact Kernel-Side Frontier.)}
The first D1-side arrow is reduced to proving four uniform sub-laws:
$\Xi \to \psi_\rho$; $\Xi \to \mathcal{O}_\rho$;
$\Xi \to \mathcal{Q}_\rho$; $\Xi \to \mathcal{L}_\rho$.
The D1 chain then sharpens to:
$\Xi \to (\psi_\rho, \mathcal{O}_\rho, \mathcal{Q}_\rho,
\mathcal{L}_\rho) \to A_\rho \to T_\rho$.
\end{corollary}

\begin{proof}
Direct consequence of Theorem~\ref{thm:rh-d1-6}. \qed
\end{proof}

\begin{lemma}[\status{C}]\label{lem:rh-d1-7}
\textup{(RH.D1.7: Orbit Extraction Minimality.)}
In the kernel package $\mathcal{K}_\rho$, the datum
$\mathcal{O}_\rho$ is the first irreducibly arithmetic ingredient:
$\psi_\rho$ alone gives only Mellin-side analytic seed data;
$\mathcal{O}_\rho$ is exactly the ingredient that turns seed data
into arithmetic-scaling data; without $\mathcal{O}_\rho$, the
descent cannot lawfully land in $A_\rho$.
\end{lemma}

\begin{proof}
\textbf{Step 1.}  $\psi_\rho$ provides Mellin-side seed structure:
a Mellin transform says the object lives in the right analytic
language, but specifies no arithmetic scaling law or orbit
organization.  So $\psi_\rho$ alone cannot determine $A_\rho$.

\textbf{Step 2.}  The RH branch descends to an arithmetic-scaling
support skeleton (not an arbitrary support object), and SCC-M2
forces admissible probes to respect multiplicative scaling
$x \mapsto x/n$.  Some kernel ingredient must encode that orbit
structure before support extraction can begin.

\textbf{Step 3.}  $\mathcal{O}_\rho$ is by definition the
arithmetic orbit/scaling datum --- the piece that groups support
components into scaling packets.  Without it, the remaining data
$(\psi_\rho, \mathcal{Q}_\rho, \mathcal{L}_\rho)$ describes a
lawful Mellin object in a quotient/Hilbert sense, but does not
specify an arithmetic-scaling subconfiguration.  The descent has no
lawful way to land in $A_\rho$. \qed
\end{proof}

\begin{corollary}[\status{C}]\label{cor:rh-d1-7-1}
\textup{(RH.D1.7.1: D1 First Arithmetic Hinge.)}
The first genuinely arithmetic hinge in the D1 chain
$\Xi \to (\psi_\rho, \mathcal{O}_\rho, \mathcal{Q}_\rho,
\mathcal{L}_\rho) \to A_\rho \to T_\rho$
is the extraction of $\mathcal{O}_\rho$. The D1-side question
sharpens to: what is the weakest uniform theorem recovering the
arithmetic orbit/scaling datum $\mathcal{O}_\rho$ from the
Mellin/Dirichlet side of $\Xi$?
\end{corollary}

\begin{proof}
By Lemma~\ref{lem:rh-d1-7}, $\mathcal{O}_\rho$ is the first
irreducibly arithmetic ingredient; the others are analytic
infrastructure.  The named subquestion is exactly the content of
the remaining D1 obligation at the orbit level. \qed
\end{proof}

\begin{theorem}[\status{C}]\label{thm:rh-d1-8}
\textup{(RH.D1.8: Orbit-Side Iff Schema.)}
The orbit-extraction sublaw $\Xi \to \mathcal{O}_\rho$ exists
lawfully if and only if there exists one uniform arithmetic-scaling
orbit grammar on the Mellin/Dirichlet side of $\Xi$ whose instances
attach to every nontrivial zero $\rho$.
\end{theorem}

\begin{proof}
\textbf{($\Leftarrow$) Sufficiency.}
Assume one uniform orbit grammar exists on $\Xi$.  Applying it to
each $\rho$ yields $\mathcal{O}_\rho$ with the same arithmetic-scaling
structure.  Since the grammar is already arithmetic-scaling by
construction, $\mathcal{O}_\rho$ lawfully seeds $A_\rho$.

\textbf{($\Rightarrow$) Necessity.}
Assume the orbit sublaw exists lawfully for every $\rho$.  By
Lemma~\ref{lem:rh-d1-7}, $\mathcal{O}_\rho$ is the first irreducibly
arithmetic ingredient.  If the rule changed from zero to zero, D1
would still depend on undeclared zero-specific choices --- violating
Book~VII's anti-smuggling rule.  Hence one common grammar exists. \qed
\end{proof}

\begin{corollary}[\status{C}]\label{cor:rh-d1-8-1}
\textup{(RH.D1.8.1: Exact First Arithmetic Subfrontier.)}
The first truly arithmetic subfrontier of D1 is: prove one uniform
arithmetic-scaling orbit grammar on the Mellin/Dirichlet side of $\Xi$.
\end{corollary}

\begin{proof}
Direct consequence of Theorem~\ref{thm:rh-d1-8}. \qed
\end{proof}

%% ---------------------------------------------------------------
%% ---------------------------------------------------------------
\section{Uniform Arithmetic-Scaling Orbit Grammar (RH-OG Series)}
\label{sec:rh-og}
%% ---------------------------------------------------------------

This section formalizes the first arithmetic subfrontier identified in
Cor.~\ref{cor:rh-d1-8-1}: proving one uniform arithmetic-scaling orbit
grammar on the Mellin/Dirichlet side of $\Xi$.

\subsection{Orbit Grammar Definition and Separation}

\begin{definition}[\status{D}]\label{def:rh-orbit-grammar}
\textup{(Uniform Arithmetic-Scaling Orbit Grammar on $\Xi$.)}
A \emph{uniform arithmetic-scaling orbit grammar} on the
Mellin/Dirichlet side of $\Xi$ is a single lawful rule system
\[
  G_{\mathrm{orb}}
  \;=\;
  (\mathrm{Obj},\,\mathrm{Prim},\,\mathrm{Rep},\,\mathrm{Scale},\,
   \mathrm{Wt},\,\mathrm{Quot},\,\mathrm{Unif})
\]
with the following data and properties:
\begin{enumerate}[label=\textup{(\arabic*)}]
  \item \emph{Arithmetic object class.}  $\mathrm{Obj}$ assigns a
    branch-visible class of Mellin/Dirichlet-side arithmetic objects
    from which local orbit/scaling data may be extracted.
  \item \emph{Primitive-orbit selector.}  $\mathrm{Prim}$ selects
    the primitive members of $\mathrm{Obj}$: the basic arithmetic
    orbit carriers.
  \item \emph{Repetition law.}  $\mathrm{Rep}$ assigns to each
    primitive object its lawful repetitions, so that iterated
    arithmetic contribution is represented without introducing new
    primitive types.
  \item \emph{Scaling law.}  $\mathrm{Scale}$ assigns to each
    primitive and its repetitions a multiplicative scaling parameter
    compatible with the RH probe discipline, uniformly in $\rho$,
    compatible with the admissible multiplicative scaling constraint.
  \item \emph{Weight law.}  $\mathrm{Wt}$ assigns the
    Mellin/Dirichlet-side weight carried by each primitive and its
    repetitions.
  \item \emph{Quotient-readiness.}  $\mathrm{Quot}$ certifies that
    the extracted orbit/scaling data is already compatible with the
    declared $\mathbb{Q}^\times$-descent side of the kernel package
    and does not depend on undeclared quotient choices.
  \item \emph{Uniformity / no-smuggling.}  $\mathrm{Unif}$ requires
    that the same rule family governs all nontrivial zeros $\rho$;
    no stage of the extraction may depend on zero-specific
    supplementation, local patching, or presentation-level choices
    not already licensed by the branch imports.
\end{enumerate}
Given $G_{\mathrm{orb}}$, the associated orbit datum $\mathcal{O}_\rho$
is extracted at probe $\rho$ by applying the common rule family to the
Mellin/Dirichlet side of $\Xi$.
\end{definition}

\begin{remark}[\status{R}]\label{rem:rh-og-role}
The role of $G_{\mathrm{orb}}$ is not to replace the full kernel
package $\mathcal{K}_\rho = (\psi_\rho, \mathcal{O}_\rho,
\mathcal{Q}_\rho, \mathcal{L}_\rho)$, but to isolate the precise
point at which Mellin-side analytic seed data becomes
arithmetic-scaling data.  This is exactly the hinge identified in
Lem.~\ref{lem:rh-d1-7} and Cor.~\ref{cor:rh-d1-7-1}.
\end{remark}

\begin{lemma}[\status{C}]\label{lem:rh-og-separation}
\textup{(Orbit Grammar Strictly Stronger Than Mellin Seed Data.)}
Within the D1 kernel package, the existence of Mellin-seed data
$\psi_\rho$ does not by itself determine a lawful arithmetic
orbit/scaling datum $\mathcal{O}_\rho$.  Unless a uniform
arithmetic-scaling orbit grammar exists, the Mellin-side seed data
remains analytic only and does not lawfully specify the arithmetic
packet structure required for descent to $A_\rho$.
\end{lemma}

\begin{proof}
Restatement of the minimality mechanism in Lem.~\ref{lem:rh-d1-7}.
$\psi_\rho$ gives only Mellin-side analytic seed structure, while
$\mathcal{O}_\rho$ is the datum that groups support components into
scaling packets.  Without such grouping, the remaining data
$(\psi_\rho, \mathcal{Q}_\rho, \mathcal{L}_\rho)$ may define a
lawful Mellin object in quotient/Hilbert terms, but not an
arithmetic-scaling subconfiguration.  Hence the descent cannot
lawfully land in $A_\rho$. \qed
\end{proof}

\begin{proposition}[\status{C}]\label{prop:rh-og-candidate-filter}
\textup{(Minimal Admissibility Filter for Candidate Orbit Grammars.)}
Any candidate realization of $G_{\mathrm{orb}}$ strong enough to
discharge the first arithmetic subfrontier must satisfy, at minimum:
\begin{enumerate}[label=\textup{(\arabic*)}]
  \item \emph{Prime-power repetition compatibility.}  The repetition
    law must encode iterative arithmetic contribution compatible with
    prime-power layering.
  \item \emph{Explicit-formula compatibility.}  The weight and
    repetition laws must be capable of reproducing the prime-power
    side of the explicit-formula structure.  This is the single most
    constraining filter.
  \item \emph{Determinant/symmetry compatibility.}  The grammar must
    not conflict with the even symmetry and growth behavior of $\Xi$.
  \item \emph{Uniformity compatibility.}  The same grammar must apply
    at every $\rho$ with no zero-specific adjustments.
\end{enumerate}
\end{proposition}

\begin{proof}
These are the weakest admissibility checks forced jointly by the
current RH frontier shape and the speculative screening already
recorded for viable RH operator/orbit architectures (SCC Branch,
Def.~\ref{def:RH-necessary-conditions}).  Explicit-formula
compatibility (item~2) is the most constraining because the
identification ``primes $\leftrightarrow$ primitive periodic orbits''
is the central structural match in the surviving candidate direction.
A grammar failing any of (1)--(4) cannot serve as the common law
behind $\Xi \to \mathcal{O}_\rho$. \qed
\end{proof}

\subsection{Scaling-Flow Candidate Architecture}

\begin{definition}[\status{D}]\label{def:rh-og-scaling-candidate}
\textup{(Scaling-Flow Orbit-Grammar Candidate $G_{\mathrm{sf}}$.)}
A \emph{scaling-flow candidate realization} of the uniform
arithmetic-scaling orbit grammar is a candidate rule system
\[
  G_{\mathrm{sf}}
  \;=\;
  (X_{\mathrm{sf}},\,\phi_t,\,\Pi_{\mathbb{Q}^\times},\,
   \mathrm{Prim},\,\mathrm{Rep},\,\mathrm{Scale},\,\mathrm{Wt})
\]
consisting of:
\begin{enumerate}[label=\textup{(\arabic*)}]
  \item a branch-visible multiplicative state space $X_{\mathrm{sf}}$;
  \item a scaling action $\phi_t(x) = e^t x$;
  \item a declared arithmetic quotient map
    $\Pi_{\mathbb{Q}^\times}: X_{\mathrm{sf}} \to
    X_{\mathrm{sf}}/\mathbb{Q}^\times$;
  \item a primitive-orbit selector $\mathrm{Prim}$ on closed orbits
    of the quotient flow;
  \item a repetition law $\mathrm{Rep}$ assigning iterates of
    primitive closed orbits;
  \item a scaling law $\mathrm{Scale}$ assigning period
    $T(\gamma_p) = \log p$ to the primitive orbit corresponding to
    prime $p$, and $T(\gamma_p^m) = m\log p$ to its $m$-fold
    repetition;
  \item a weight law $\mathrm{Wt}$ assigning the half-density orbit
    weight $\mathrm{Wt}(\gamma_p^m) = e^{-T(\gamma_p^m)/2} =
    p^{-m/2}$.
\end{enumerate}
A candidate $G_{\mathrm{sf}}$ is RH-admissible at the orbit level
only if all assignments above are uniform in $\rho$ and require no
zero-specific supplementation.
\end{definition}

\begin{remark}[\status{R}]\label{rem:rh-sf-candidate-status}
This definition records a \emph{candidate} instance of the RH.D1.8
orbit grammar.  It is not theorem-bearing discharge.  Its role is to
give one concrete architecture against which the canonical
admissibility tests may be applied.  The scaling/adelic flow picture
and transfer-operator route to $\zeta(s)$ are not NFC-original; any
canonical use must explicitly attribute the external literature and
claim only the force licensed by proved NFC branch-level admissibility
theorems.
\end{remark}

\begin{theorem}[\status{C}]\label{thm:rh-og-1}
\textup{(Conditional Scaling-Flow Realization of the Orbit Grammar.)}
\textup{[Conditional on: prime/orbit bijection; repetition law;
half-density weight; and uniform quotient-visibility of all
assignments.]}
If these four conditions hold for $G_{\mathrm{sf}}$, then
$G_{\mathrm{sf}}$ defines a uniform arithmetic-scaling orbit grammar in
the sense of Def.~\ref{def:rh-orbit-grammar}, and therefore the
orbit-extraction sublaw $\Xi \to \mathcal{O}_\rho$ is conditionally
realized by the scaling-flow architecture.
\end{theorem}

\begin{proof}
Conditions (i)--(ii) supply the primitive-orbit and repetition laws
required by Def.~\ref{def:rh-orbit-grammar}.  Condition~(iii) gives
the Mellin-compatible weighting law.  Condition~(iv) enforces the
Book~VII anti-smuggling requirement: the same rule family applies for
every $\rho$, with no local patching.  Hence the data of $G_{\mathrm{sf}}$
instantiate a uniform arithmetic-scaling orbit grammar.
By Thm.~\ref{thm:rh-d1-8} the lawful orbit-extraction sublaw follows
conditionally. \qed
\end{proof}

\begin{proposition}[\status{C}]\label{prop:rh-og-2}
\textup{(Scaling-Flow Admissibility Reduces to Two Checks.)}
For the scaling-flow candidate $G_{\mathrm{sf}}$, discharge of the
first arithmetic subfrontier reduces to verification of:
\begin{enumerate}[label=\textup{(\arabic*)}]
  \item \emph{Explicit-formula compatibility.}  The associated
    trace law reproduces the prime-power terms
    $\sum_{p,m} (\log p)\,p^{-m/2}\,h(m\log p)$ required by
    RH-N2.
  \item \emph{No-smuggling compatibility.}  The quotient flow,
    orbit assignments, and weight law determine $\mathcal{O}_\rho$
    by one common rule family, with no undeclared auxiliary arithmetic
    input and no zero-specific choices.
\end{enumerate}
If both checks hold, $G_{\mathrm{sf}}$ is a lawful candidate discharge
of the first arithmetic subfrontier (Cor.~\ref{cor:rh-d1-8-1}).
\end{proposition}

\begin{proof}
By the SCC/RH screening package
(Prop.~\ref{prop:RH-architecture-elimination}, SCC Branch), RH-N2 is
the single most constraining pre-construction filter.  The
identification ``primes $\leftrightarrow$ primitive periodic orbits''
is the central structural match.  A proposed grammar failing check~(1)
does not satisfy the main admissibility gate; failing check~(2) violates
Book~VII no-smuggling discipline. \qed
\end{proof}

\subsection{Scaling-Flow Period Law and Half-Density Weight}

\begin{definition}[\status{D}]\label{def:rh-trace-template}
\textup{(Branch-Admissible RH Trace Template.)}
A \emph{branch-admissible RH trace template} for a candidate orbit
grammar $G$ is a rule assigning, to every compactly supported even
test function $h$, a distributional trace expression
$\mathrm{Tr}_G(h)$ such that:
\begin{enumerate}[label=\textup{(\arabic*)}]
  \item primitive periodic orbits contribute by a uniform common law;
  \item repetitions contribute by the iterate law;
  \item the contribution attached to a primitive orbit of prime label
    $p$ and repetition $m$ is proportional to
    $(\log p)\,p^{-m/2}\,h(m\log p)$;
  \item the same rule applies for every zero parameter $\rho$;
  \item no undeclared auxiliary arithmetic data enters the
    construction.
\end{enumerate}
\end{definition}

\begin{proposition}[\status{C}]\label{prop:rh-sf-period-law}
\textup{(Scaling-Flow Period Law.)}
For the scaling-flow candidate $G_{\mathrm{sf}}$ with primitive closed
orbits in bijection with primes, the lawful period assignment is
\[
  T(\gamma_p) = \log p,\qquad
  T(\gamma_p^m) = m\log p.
\]
\end{proposition}

\begin{proof}
In $G_{\mathrm{sf}}$, closed orbit closure is the condition that the
scaling factor returns to the same arithmetic class modulo
$\mathbb{Q}^\times$.  By the candidate architecture
(Def.~\ref{def:rh-og-scaling-candidate}), primitive closure occurs at
$e^T = p$, giving $T = \log p$ for a prime-labeled primitive orbit.
Iteration yields $m\log p$ for the $m$-fold repetition.  This is
precisely the prime-power repetition law required by the orbit grammar. \qed
\end{proof}

\begin{lemma}[\status{C}]\label{lem:rh-sf-half-density-weight}
\textup{(Half-Density Orbit Weight.)}
Under $G_{\mathrm{sf}}$ with half-density normalization $\alpha =
\tfrac{1}{2}$, the canonical orbit weight on the $m$-fold iterate of
$\gamma_p$ is
\[
  \mathrm{Wt}(\gamma_p^m)
  \;=\; e^{-\alpha\,T(\gamma_p^m)}
  \;=\; p^{-m/2}.
\]
\end{lemma}

\begin{proof}
By Prop.~\ref{prop:rh-sf-period-law}, $T(\gamma_p^m) = m\log p$.
Substituting into half-density normalization:
$e^{-(1/2)m\log p} = p^{-m/2}$.  This matches the prime-power
damping required by the explicit-formula side. \qed
\end{proof}

\subsection{Determinant Channel and Log-Derivative Route}

\begin{proposition}[\status{C}]\label{prop:rh-sf-det-euler}
\textup{(Determinant Channel for the Scaling-Flow Candidate.)}
\textup{[Conditional on branch-visibility of the transfer-operator
family $(\mathcal{S}_s f)(x) = \sum_{n\geq 1} n^{-s} f(x/n)$ and
its admissible determinant channel.]}
Within its declared scope,
\[
  \det(1 - \mathcal{S}_s)^{-1}
  \;=\; \prod_p (1-p^{-s})^{-1}
  \;=\; \zeta(s).
\]
The scaling-flow candidate thereby has a lawful determinant channel
compatible with the orbit grammar target for $\mathcal{O}_\rho$.
\end{proposition}

\begin{proof}
Under the branch-visibility hypothesis, the resulting determinant
identity belongs to declared RH branch data
(via $W_{RH}$, $T_{RH}$) rather than to undeclared external
supplementation.  This is the determinant-level hardening of the
candidate pipeline
$\text{scaling flow} \to \mathcal{S}_s \to
\det(1-\mathcal{S}_s)^{-1} = \zeta(s)$. \qed
\end{proof}

\begin{lemma}[\status{C}]\label{lem:rh-sf-primepower-logderiv}
\textup{(Prime-Power Logarithmic-Derivative Law.)}
\textup{[Conditional on Prop.~\ref{prop:rh-sf-det-euler}.]}
The differentiated determinant channel has formal prime-power
expansion
\[
  -\frac{d}{ds}\log\det(1-\mathcal{S}_s)^{-1}
  \;=\; -\frac{\zeta'(s)}{\zeta(s)}
  \;=\; \sum_p \sum_{m\geq 1} (\log p)\,p^{-ms}.
\]
\end{lemma}

\begin{proof}
By Prop.~\ref{prop:rh-sf-det-euler}: $\log\det(1-\mathcal{S}_s)^{-1}
= \sum_p \sum_{m\geq 1} p^{-ms}/m$.  Differentiating with respect to
$s$: differentiating $p^{-ms}$ gives $-m\log p$, which cancels $1/m$
from the logarithm expansion, yielding $\sum_{p,m} (\log p)\,p^{-ms}$.
The $\log p$ factor is therefore the logarithmic-derivative weight of
the Euler-product channel --- not an added normalization. \qed
\end{proof}

\begin{corollary}[\status{C}]\label{cor:rh-sf-logp-template}
\textup{(RH-N2 Trace Template from the Determinant Channel.)}
\textup{[Conditional on Lem.~\ref{lem:rh-sf-primepower-logderiv},
Prop.~\ref{prop:rh-sf-period-law}, Lem.~\ref{lem:rh-sf-half-density-weight}.]}
At $s = \tfrac{1}{2}+it$, the prime-power part of the admissible
trace pairing against any even test function $h$ has the required
RH-N2 form:
\[
  \sum_{p,m} (\log p)\,p^{-m/2}\,h(m\log p).
\]
\end{corollary}

\begin{proof}
The period law supplies the orbit times $m\log p$; the half-density
normalization supplies $p^{-m/2}$; the log-derivative channel
contributes the exact coefficient $\log p$.  Pairing with an even
test function yields the RH-N2 prime-power template. \qed
\end{proof}

\subsection{Open Obligations in the Orbit Grammar Workstream}

\begin{openobligation}[\status{C}]\label{ob:rh-sf-logderiv-legality}
\textup{(Log-Derivative Legality in Branch Scope.)}
\textup{[Conditionally discharged by Def.~\ref{def:rh-logdiff-channel}
together with Props.~\ref{prop:rh-sf-operator-origin}
and~\ref{prop:rh-sf-determinant-origin}.]}
The determinant/log-derivative manipulations are lawful on the
declared RH witness family: $D_{RH}(s)$ is a declared branch object
(Prop.~\ref{prop:rh-sf-determinant-origin}); $D \mapsto
-\tfrac{d}{ds}\log D$ is a licensed branch operation on
$\mathcal{O}_{RH}^{\det}$ (Def.~\ref{def:rh-logdiff-channel}
Clause~2); interface lawfulness holds
(Rem.~\ref{rem:rh-logdiff-interface-lawfulness}); no undeclared
structure is imported
(Rem.~\ref{rem:rh-logdiff-attribution-discipline}).
All four clauses of Def.~\ref{def:rh-logdiff-channel} are satisfied.
\end{openobligation}

\begin{openobligation}[\status{C}]\label{ob:rh-sf-trace-pairing-law}
\textup{(Admissible Trace Pairing Law.)}
\textup{[Conditionally discharged by
Prop.~\ref{prop:rh-sf-primepower-pairing} and
Lem.~\ref{lem:rh-sf-trace-uniformity}.]}
For every $h \in H_{RH}^{\mathrm{tr}}$ (Def.~\ref{def:rh-trace-test-subfamily}),
the differentiated determinant channel pairs with declared RH
trace-test objects in $W_{RH}$ to produce
$\sum_{p,m}(\log p)\,p^{-m/2}\,h(m\log p)$
(Prop.~\ref{prop:rh-sf-primepower-pairing}),
and this realization is branch-visible and uniform in $\rho$
(Lem.~\ref{lem:rh-sf-trace-uniformity}).
The RH-N2 prime-power trace template is realized in
branch-admissible form.
\end{openobligation}

\begin{theorem}[\status{C}]\label{thm:rh-og-5}
\textup{(Conditional RH-N2 Realization via the Log-Derivative Channel.)}
\textup{[Conditional on: Prop.~\ref{prop:rh-sf-det-euler};
Prop.~\ref{prop:rh-sf-period-law};
Lem.~\ref{lem:rh-sf-half-density-weight};
ob:rh-sf-logderiv-legality discharged;
ob:rh-sf-trace-pairing-law discharged.]}
The scaling-flow candidate satisfies RH-N2 in branch-admissible
form.  Consequently, it passes the main explicit-formula gate for a
conditional realization of the first arithmetic subfrontier underlying
Thm.~\ref{thm:rh-d1-8}.
\end{theorem}

\begin{proof}
By Prop.~\ref{prop:rh-sf-det-euler}: lawful determinant channel.  By
Props.~\ref{prop:rh-sf-period-law} and~\ref{lem:rh-sf-half-density-weight}:
correct orbit lengths and half-density damping.  By
ob:rh-sf-logderiv-legality: the $\log p$ amplitude is lawful in
branch scope.  By ob:rh-sf-trace-pairing-law: the differentiated
channel pairs with declared RH trace-test objects to realize the
prime-power explicit-formula data as periodic-orbit data.  This is
exactly RH-N2.  Since RH-N2 is the single most constraining filter
(Prop.~\ref{prop:RH-architecture-elimination}, SCC Branch), the
candidate clears the main explicit-formula admissibility gate. \qed
\end{proof}

\subsection{D2-Admissibility Checklist and Reduction}

\begin{definition}[\status{D}]\label{def:rh-d2-checklist}
\textup{(D2-Admissibility Checklist.)}
A candidate spectral/trace construction for the RH branch is
\emph{D2-admissible} if and only if all of the following hold:
\begin{enumerate}[label=\textup{(\arabic*)}]
  \item \emph{Declared object law.}  Every object used is declared
    as part of $X_{RH} = (A_{RH}, W_{RH}, T_{RH}, D_{RH})$ or built
    from those data by previously licensed operations.
  \item \emph{No external operator import.}  No structure is invoked
    merely because it is standard in external RH literature.
  \item \emph{Witness-family visibility.}  Every test object belongs
    to the declared witness family $W_{RH}$.
  \item \emph{Transported-invariant visibility.}  Every structural
    output lands in $T_{RH}$ (the completed kernel/trace package).
  \item \emph{Uniformity in $\rho$.}  The same rule family applies
    for every nontrivial zero $\rho$; no zero-by-zero adjustment
    activates anti-smuggling failure.
  \item \emph{Interface lawfulness.}  Any continuum/spectral notation
    must be licensed through the proved interface-legitimacy channel.
  \item \emph{Attribution without force inflation.}  Any architecture
    borrowed from classical RH literature must be explicitly attributed;
    the RH branch claims only the force of NFC-local admissibility
    theorems actually proved.
\end{enumerate}
\end{definition}

\begin{openobligation}[\status{C}]\label{ob:rh-sf-d2-audit}
\textup{(D2 Audit for the Scaling-Flow Determinant Channel.)}
Prove that the scaling-flow/transfer-operator/determinant/log-derivative
construction satisfies Def.~\ref{def:rh-d2-checklist}.  Specifically:
(1) $\mathcal{S}_s$ is branch-visible in the declared RH witness/probe
regime; (2) $\det(1-\mathcal{S}_s)^{-1}$ is a lawful branch object;
(3) the differentiated channel pairing is realized by declared
trace-test objects in $W_{RH}$; (4) the resulting prime-power trace
distribution lands in $T_{RH}$; (5) the whole construction is uniform
in $\rho$.
\end{openobligation}

\begin{lemma}[\status{C}]\label{lem:rh-d2-reduction}
\textup{(D2 Reduces to Four Local Certification Checks.)}
For the scaling-flow candidate, RH-AC-D2 is discharged if four local
certification checks are proved:
\begin{enumerate}[label=\textup{(\arabic*)}]
  \item \emph{Operator-origin check.}  $\mathcal{S}_s$ is
    reconstructed from declared Mellin/scaling probe data in $W_{RH}$,
    not imported as an external black box.
  \item \emph{Determinant-origin check.}  The
    determinant/log-derivative channel is defined by a previously
    licensed branch operation on that family.
  \item \emph{Trace-test check.}  The pairing producing the
    prime-power distribution is performed against declared RH
    trace-test objects in $W_{RH}$.
  \item \emph{Package-landing check.}  The outcome is recorded inside
    $T_{RH} = (\Xi, K_{RH}, \mathcal{T}_{pp}, \mathcal{Q}_{RH})$.
\end{enumerate}
\end{lemma}

\begin{proof}
RH-AC-D2 is the failure mode in which a purported branch-visible datum
depends on undeclared external analytic/operator structure.  The RH
object tells us exactly where lawful data may live: $W_{RH}$ for
probes/test objects and $T_{RH}$ for transported invariant output.
A D2 discharge is therefore nothing more and nothing less than proving
that operator, determinant, trace pairing, and final output all stay
inside those declared channels. \qed
\end{proof}

\begin{proposition}[\status{C}]\label{prop:rh-sf-logderiv-d2-ready}
\textup{(Readiness Criterion for the Log-Derivative Route.)}
Assuming Prop.~\ref{prop:rh-sf-det-euler},
Lem.~\ref{lem:rh-sf-primepower-logderiv}, and
Lem.~\ref{lem:rh-d2-reduction}, the only remaining barrier to using
the scaling-flow log-derivative channel as a lawful RH-N2 witness is
the proof of the four local certification checks in
Lem.~\ref{lem:rh-d2-reduction}.
\end{proposition}

\begin{proof}
The determinant/log-derivative packet already explains the formal
origin of the $\log p$ coefficient.  By Lem.~\ref{lem:rh-d2-reduction},
the D2 obstruction is equivalent to the four local certification checks.
Those checks are now the entire remaining barrier. \qed
\end{proof}

\begin{remark}[\status{R}]\label{rem:rh-sf-remaining-burden}
\textbf{Project-level compression.}
After this package, the first arithmetic subfrontier is no longer
``find an orbit grammar''.  The scaling-flow candidate $G_{\mathrm{sf}}$
has already resolved:
\begin{itemize}
  \item primitive-orbit identification (primes),
  \item repetition law (prime powers),
  \item period law ($m\log p$),
  \item half-density damping ($p^{-m/2}$),
  \item formal log-derivative route to the $\log p$ amplitude
    (Lem.~\ref{lem:rh-sf-primepower-logderiv}).
\end{itemize}
The remaining burden is exactly two items:
\begin{enumerate}[label=\textup{(\arabic*)}]
  \item Prove the log-derivative manipulations are branch-visible
    and no-smuggling (ob:rh-sf-logderiv-legality,
    ob:rh-sf-d2-audit).
  \item Prove the resulting prime-power distribution pairs correctly
    with declared RH trace-test objects (ob:rh-sf-trace-pairing-law).
\end{enumerate}
This is the exact content of the first arithmetic subfrontier
after all structural reductions.
\end{remark}


%% ---------------------------------------------------------------
\section{First Hard Block: D2 Audit, Log-Derivative, Trace Pairing}
\label{sec:rh-first-hard-block}
%% ---------------------------------------------------------------

This section develops the complete closing ladder for the first RH
hard block.  The mission: prove the scaling-flow determinant and
log-derivative route is entirely branch-internal, pairs with declared
RH trace-test objects, and lands in the declared transported package,
uniformly in $\rho$.

\subsection{Scaling-Flow Probe Subfamily and Reconstruction}

\begin{definition}[\status{D}]\label{def:rh-sf-probe-subfamily}
\textup{(Scaling-Flow Probe Subfamily $W_{RH}^{\mathrm{sf}}$.)}
Let $W_{RH}^{\mathrm{sf}} \subseteq W_{RH}$ be the smallest declared
subfamily satisfying:
\begin{enumerate}[label=\textup{(\arabic*)}]
  \item it contains the Mellin/scaling probe generators arising from
    the declared Mellin/Dirichlet descent data;
  \item it is closed under the branch-visible arithmetic-scaling
    action $f(x) \mapsto f(x/n)$ for every admissible scaling
    index $n$;
  \item it is closed under the branch-visible Mellin weighting rule
    attached to the descent parameter $s$;
  \item it is closed under the licensed finite/admissibly controlled
    aggregation rule used to combine scaled probes.
\end{enumerate}
This is the smallest declaration that makes ``$\mathcal{S}_s$ comes
from $W_{RH}$'' a precise branch statement.
\end{definition}

\begin{definition}[\status{D}]\label{def:rh-sf-aggregation-channel}
\textup{(Scaling-Stable Mellin Aggregation Operation.)}
Let
\[
  \mathrm{Agg}^{\mathrm{sf}}_s : \mathrm{Fam}(W_{RH}^{\mathrm{sf}})
  \to W_{RH}^{\mathrm{sf}}
\]
be the scaling-stable Mellin aggregation operation defined by
\[
  \mathrm{Agg}^{\mathrm{sf}}_s\bigl(\{f_n\}_n\bigr)(x)
  \;:=\; \sum_{n \in \mathrm{AdmScale}} n^{-s} f_n(x),
\]
whenever $\{f_n\}_n$ is generated from the declared arithmetic-scaling
descent channel, each $f_n \in W_{RH}^{\mathrm{sf}}$, the index set
$\mathrm{AdmScale}$ is branch-declared, the weight $n^{-s}$ belongs
to the Mellin/Dirichlet descent channel, and the aggregation stays
inside $W_{RH}^{\mathrm{sf}}$.
\end{definition}

\begin{definition}[\status{D}]\label{def:rh-sf-reconstruction-rule}
\textup{(Scaling-Flow Reconstruction Operator $\mathcal{R}_s$.)}
Define
\[
  \mathcal{R}_s(f)
  \;:=\; \mathrm{Agg}^{\mathrm{sf}}_s
  \bigl(\{x \mapsto f(x/n)\}_{n \in \mathrm{AdmScale}}\bigr).
\]
When context is clear, write $\mathcal{S}_s := \mathcal{R}_s$.
The transfer operator $\mathcal{S}_s$ is not a primitive; it is
\emph{reconstructed} from two declared local ingredients: scaling of
witness probes, and Mellin-weighted aggregation of those scaled probes.
\end{definition}

\begin{proposition}[\status{C}]\label{prop:rh-sf-operator-origin}
\textup{(Operator-Origin of the Scaling-Flow Channel.)}
\textup{[Conditional on declared Mellin/Dirichlet descent data and
SCC-M1--M4 admissibility.]}
The scaling-flow operator family $(\mathcal{S}_s f)(x) = \sum_n
n^{-s} f(x/n)$ is reconstructed from declared Mellin/scaling probe
data in $W_{RH}$, and is therefore branch-visible rather than
externally imported.  In particular, the operator-origin check of
Lem.~\ref{lem:rh-d2-reduction} is satisfied.
\end{proposition}

\begin{proof}
By Def.~\ref{def:rh-sf-probe-subfamily}, $W_{RH}^{\mathrm{sf}}$ is
a declared branch subfamily.  For $f \in W_{RH}^{\mathrm{sf}}$, the
scaling action $x \mapsto f(x/n)$ is a licensed endomorphism on the
probe class (item~2 of the definition).  The factor $n^{-s}$ is the
Mellin weight of the same declared descent channel, not an added
analytic decoration.  The sum is the licensed aggregation operation
$\mathrm{Agg}^{\mathrm{sf}}_s$ (Def.~\ref{def:rh-sf-aggregation-channel}).
Hence $\mathcal{S}_s = \mathcal{R}_s$ is reconstructed from declared
probe data rather than imported as an external black box, satisfying
the operator-origin check. \qed
\end{proof}

\subsection{Determinant Channel}

\begin{definition}[\status{D}]\label{def:rh-det-channel-local}
\textup{(Determinant Channel $\mathrm{Det}_{RH}^{\mathrm{sf}}$.)}
Let $\mathcal{F}_{RH}^{\mathrm{sf}}$ be the declared family generated
by the branch-visible endomorphisms $(1-\mathcal{S}_s)$ acting on
$W_{RH}^{\mathrm{sf}}$.  A \emph{determinant channel} for the
scaling-flow family is a branch operation
\[
  \mathrm{Det}_{RH}^{\mathrm{sf}}:
  \mathcal{F}_{RH}^{\mathrm{sf}} \to \mathcal{O}_{RH}^{\det}
\]
into a declared determinant-output object class
$\mathcal{O}_{RH}^{\det}$, producing the lawful branch object
\[
  D_{RH}(s) := \mathrm{Det}_{RH}^{\mathrm{sf}}(1-\mathcal{S}_s)^{-1}.
\]
\end{definition}

\begin{proposition}[\status{C}]\label{prop:rh-sf-determinant-origin}
\textup{(Determinant-Origin of the Scaling-Flow Channel.)}
\textup{[Conditional on Prop.~\ref{prop:rh-sf-operator-origin} and
Def.~\ref{def:rh-det-channel-local}, with interface-lawful notation.]}
The scaling-flow determinant channel $D_{RH}(s) = \mathrm{Det}^{\mathrm{sf}}_{RH}
(1-\mathcal{S}_s)^{-1}$ is a lawful branch object built from
declared RH data by a previously licensed branch operation.  The
determinant-origin check of Lem.~\ref{lem:rh-d2-reduction} is
satisfied.

Further, the comparison $D_{RH}(s) = \prod_p(1-p^{-s})^{-1}$ is
cited with attribution only: it does not supply branch objecthood.
\end{proposition}

\begin{proof}
By Prop.~\ref{prop:rh-sf-operator-origin}, $(1-\mathcal{S}_s) \in
\mathcal{F}_{RH}^{\mathrm{sf}}$ is built from declared RH witness/probe
data.  By Def.~\ref{def:rh-det-channel-local}, $\mathrm{Det}_{RH}^{\mathrm{sf}}$
is a declared branch operation on $\mathcal{F}_{RH}^{\mathrm{sf}}$.
Hence $D_{RH}(s)$ is a branch-visible object built by a previously
licensed operation, satisfying the determinant-origin check.  The
Euler-product comparison has attributed comparison force only;
legality comes from the branch construction, not from the external
identity. \qed
\end{proof}

\subsection{Log-Derivative Channel}

\begin{definition}[\status{D}]\label{def:rh-logdiff-channel}
\textup{(Logarithmic-Differentiation Channel $\mathrm{LogDiff}_{RH}$.)}
Let $\mathcal{O}_{RH}^{\det}$ be the declared determinant-output
object class.  A \emph{logarithmic-differentiation channel} is a
branch operation
\[
  \mathrm{LogDiff}_{RH} :
  \mathcal{O}_{RH}^{\det} \to \mathcal{O}_{RH}^{\log}
\]
with $\mathrm{LogDiff}_{RH}(D)(s) := -\tfrac{d}{ds}\log D(s)$,
subject to four admissibility clauses:
\begin{enumerate}[label=\textup{(\arabic*)}]
  \item \emph{Declared-input condition.}  $D(s)$ is already a lawful
    branch object in $\mathcal{O}_{RH}^{\det}$.
  \item \emph{Operation-license condition.}  The composite map
    $D \mapsto -\tfrac{d}{ds}\log D$ is itself declared as a branch
    operation on $\mathcal{O}_{RH}^{\det}$, not inferred from
    external operator lore.
  \item \emph{Interface condition.}  All differential/spectral
    notation is lawful in the declared interface regime.
  \item \emph{No-smuggling condition.}  No auxiliary analytic/operator
    structure is introduced beyond what is already declared in the
    RH witness/probe regime.
\end{enumerate}
\end{definition}

\begin{remark}[\status{R}]\label{rem:rh-logdiff-interface-lawfulness}
The channel $\mathrm{LogDiff}_{RH}$ is used only in the narrow
interface-legitimate sense required by the RH D2 admissibility
checklist (Def.~\ref{def:rh-d2-checklist}).  No general analytic
determinant calculus, spectral regularization theory, or undeclared
continuum operator package is invoked.  For determinant objects
already certified in RH scope, the composite operation $D \mapsto
-\tfrac{d}{ds}\log D$ is admitted only as a licensed local channel
on the declared determinant-output class.
\end{remark}

\begin{remark}[\status{R}]\label{rem:rh-logdiff-attribution-discipline}
Any comparison of the scaling-flow determinant channel with classical
external formulas (e.g., $\det(1-\mathcal{S}_s)^{-1} = \prod_p
(1-p^{-s})^{-1}$ and its logarithmic-derivative expansion) is cited
only for \emph{attributed comparison and output identification}.
Such comparison does not create branch objecthood, does not license
the operation $D \mapsto -\tfrac{d}{ds}\log D$, and does not
strengthen theorem force beyond declared RH branch scope.  Branch
legality comes only from Defs.~\ref{def:rh-det-channel-local}
and~\ref{def:rh-logdiff-channel}.
\end{remark}

\begin{lemma}[\status{C}]\label{lem:rh-sf-logdiff-amplitude}
\textup{(Log-Derivative Amplitude Identification.)}
\textup{[Conditional on Prop.~\ref{prop:rh-sf-determinant-origin},
Def.~\ref{def:rh-logdiff-channel}, and ob:rh-sf-logderiv-legality.]}
The lawful logarithmic-differentiation channel applied to
$D_{RH}(s)$ identifies the prime-power amplitude contribution as
$\mathrm{LogDiff}_{RH}(D_{RH})(s) \leadsto \sum_{p,m} (\log p)\,
p^{-ms}$, so the coefficient attached to the $m$-fold repetition of
the prime orbit for $p$ is exactly $\log p$.  Equivalently, within
branch scope, the $\log p$ factor appearing in the RH-N2 prime-power
trace template is the amplitude image of the lawful
log-differentiation channel on the determinant object.
\end{lemma}

\begin{proof}
By Prop.~\ref{prop:rh-sf-determinant-origin}, $D_{RH}(s)$ is a lawful
branch object.  By Def.~\ref{def:rh-logdiff-channel} and
ob:rh-sf-logderiv-legality, applying $-\tfrac{d}{ds}\log$ is a
lawful branch operation.  The attributed comparison
$D_{RH}(s) \widehat{=} \prod_p(1-p^{-s})^{-1}$ is used only to
identify the output pattern: differentiating the Euler-product
comparison yields $-\tfrac{d}{ds}\log D_{RH}(s) \widehat{=}
\sum_{p,m}(\log p)\,p^{-ms}$, where $-m\log p$ from differentiating
$p^{-ms}$ cancels $1/m$ from the logarithm expansion.  Legality came
from the branch channel; the comparison only identifies the value.
Hence $\log p$ is available in branch scope as the amplitude output
of $\mathrm{LogDiff}_{RH}$. \qed
\end{proof}

\subsection{Trace-Test Subfamily and Pairing}

\begin{definition}[\status{D}]\label{def:rh-trace-test-subfamily}
\textup{(RH Trace-Test Subfamily $H_{RH}^{\mathrm{tr}}$.)}
Let $H_{RH}^{\mathrm{tr}} \subseteq W_{RH}$ be the declared RH
trace-test subfamily consisting of witness objects $h$ satisfying:
\begin{enumerate}[label=\textup{(\arabic*)}]
  \item \emph{Witness-family visibility}: $h \in W_{RH}$;
  \item \emph{Uniform admissibility}: the same rule applies for every
    $\rho$, with no zero-specific auxiliary data;
  \item \emph{Even-test symmetry}: $h$ is an even trace-test object;
  \item \emph{Prime-power receptivity}: $h$ is lawful input for the
    branch-admissible trace pairing channel;
  \item \emph{No hidden supplementation}: no extra structure beyond
    the declared RH witness/probe regime.
\end{enumerate}
\end{definition}

\begin{definition}[\status{D}]\label{def:rh-logderiv-trace-pairing}
\textup{(RH Log-Derivative Trace Pairing Map.)}
Define the RH trace-pairing map
\[
  \mathrm{Pair}_{RH}^{\mathrm{sf}} :
  \bigl(H_{RH}^{\mathrm{tr}},\, D'_{RH}\bigr)
  \longrightarrow T_{pp}
\]
by
\[
  \mathrm{Pair}_{RH}^{\mathrm{sf}}(h, D'_{RH})
  \;:=\; \sum_{p,m} (\log p)\,p^{-m/2}\,h(m\log p),
\]
whenever this sum is realized as the branch-visible pairing of the
differentiated determinant channel $D'_{RH}$ against a declared test
object $h \in H_{RH}^{\mathrm{tr}}$, with orbit times $m\log p$
from the period law, factor $p^{-m/2}$ from half-density
normalization, and coefficient $\log p$ from the log-derivative
channel.
\end{definition}

\begin{definition}[\status{D}]\label{def:rh-primepower-package}
\textup{(Prime-Power Trace Package $T_{pp}$.)}
Let $T_{pp}$ denote the prime-power trace component of the transported
RH package $T_{RH} = (\Xi, K_{RH}, T_{pp}, Q_{RH})$, where $T_{pp}$
is the declared codomain for branch-visible prime-power trace outputs
obtained from admissible trace pairing.
\end{definition}

\begin{lemma}[\status{C}]\label{lem:rh-sf-trace-uniformity}
\textup{(Uniformity of the Scaling-Flow Trace Pairing.)}
\textup{[Conditional on Defs.~\ref{def:rh-trace-test-subfamily},
\ref{def:rh-logderiv-trace-pairing} and ob:rh-sf-logderiv-legality,
with both admissibility rules stated independently of $\rho$.]}
The pairing rule $(h, \rho) \mapsto \mathrm{Pair}_{RH}^{\mathrm{sf}}
(h, D'_{RH})$ is uniform in $\rho$: the same declared admissibility
rule for $h$, the same differentiated-channel construction, and the
same pairing law apply for every nontrivial zero $\rho$, with no
zero-specific auxiliary data, normalization change, or hidden
supplementation.
\end{lemma}

\begin{proof}
By Def.~\ref{def:rh-trace-test-subfamily}, admissibility of $h$ is
fixed by one branch declaration inside $W_{RH}$, not by a
$\rho$-indexed recipe.  By ob:rh-sf-logderiv-legality,
$D'_{RH}$ is a lawful branch-visible construction, again without
$\rho$-indexed supplementation.  Since the pairing map
(Def.~\ref{def:rh-logderiv-trace-pairing}) is declared once for all,
the same rule applies for every $\rho$; the pairing is $\rho$-uniform
and introduces no hidden zero-specific supplementation. \qed
\end{proof}

\subsection{Trace-Origin and Prime-Power Pairing}

\begin{proposition}[\status{C}]\label{prop:rh-sf-trace-test-origin}
\textup{(Trace-Test Check.)}
\textup{[Conditional on Defs.~\ref{def:rh-trace-test-subfamily},
\ref{def:rh-logderiv-trace-pairing}, Props.~\ref{prop:rh-sf-operator-origin},
\ref{prop:rh-sf-determinant-origin}, lawful log-derivative channel,
and period/half-density laws.]}
The trace pairing producing the prime-power distribution is performed
against declared RH trace-test objects in $W_{RH}$.  In particular,
the trace-test local check of Lem.~\ref{lem:rh-d2-reduction} is
satisfied.
\end{proposition}

\begin{proof}
By Def.~\ref{def:rh-trace-test-subfamily}, every admissible $h$ used
in pairing lies in $W_{RH}$ with no hidden external supplementation.
By Props.~\ref{prop:rh-sf-operator-origin}
and~\ref{prop:rh-sf-determinant-origin}, the differentiated
determinant channel is a branch-visible object.  The log-derivative
channel contributes the amplitude $\log p$ (Lem.~\ref{lem:rh-sf-logdiff-amplitude});
the period and half-density laws contribute $m\log p$ and $p^{-m/2}$.
Pairing against $h \in H_{RH}^{\mathrm{tr}}$ is therefore pairing
against declared RH trace-test objects, satisfying the trace-test
check. \qed
\end{proof}

\begin{proposition}[\status{C}]\label{prop:rh-sf-primepower-pairing}
\textup{(Prime-Power Pairing Law.)}
\textup{[Conditional on Prop.~\ref{prop:rh-sf-trace-test-origin}.]}
For every $h \in H_{RH}^{\mathrm{tr}}$,
\[
  \mathrm{Pair}_{RH}^{\mathrm{sf}}(h, D'_{RH})
  \;=\; \sum_{p,m} (\log p)\,p^{-m/2}\,h(m\log p).
\]
The differentiated determinant channel realizes the RH-N2 prime-power
trace distribution as periodic-orbit data for the branch's own
admissible trace tests.
\end{proposition}

\begin{proof}
The three structural ingredients are branch-visible outputs of
previously certified channels: orbit times $m\log p$
(Prop.~\ref{prop:rh-sf-period-law}), half-density damping
$p^{-m/2}$ (Lem.~\ref{lem:rh-sf-half-density-weight}), and
log-derivative amplitude $\log p$
(Lem.~\ref{lem:rh-sf-logdiff-amplitude}).  Applying the declared
pairing map (Def.~\ref{def:rh-logderiv-trace-pairing}) to $h \in
H_{RH}^{\mathrm{tr}}$ yields exactly the stated formula.  Since the
same rule holds for every $\rho$ (Lem.~\ref{lem:rh-sf-trace-uniformity}),
the output is branch-visible and $\rho$-uniform. \qed
\end{proof}

\begin{proposition}[\status{C}]\label{prop:rh-sf-package-landing}
\textup{(Package-Landing Check.)}
\textup{[Conditional on Def.~\ref{def:rh-primepower-package} and
Prop.~\ref{prop:rh-sf-primepower-pairing}.]}
The prime-power distribution produced by admissible trace pairing
lands in the declared transported package:
\[
  \mathrm{Pair}_{RH}^{\mathrm{sf}}(h, D'_{RH})
  \;\in\; T_{pp} \subset T_{RH}
  = (\Xi, K_{RH}, T_{pp}, Q_{RH}).
\]
The package-landing local check of Lem.~\ref{lem:rh-d2-reduction} is
satisfied.
\end{proposition}

\begin{proof}
By Def.~\ref{def:rh-primepower-package}, $T_{pp}$ is the declared
codomain for branch-visible prime-power trace outputs.  By
Prop.~\ref{prop:rh-sf-primepower-pairing}, the pairing produces
exactly such an output.  Hence the output lands in $T_{pp} \subset
T_{RH}$, satisfying the package-landing check. \qed
\end{proof}

\subsection{Local D2 Audit Theorem and Closure}

\begin{theorem}[\status{C}]\label{thm:rh-sf-local-d2-checks}
\textup{(Local Certification Sufficiency for the Scaling-Flow D2 Audit.)}
\textup{[Conditional on Props.~\ref{prop:rh-sf-operator-origin},
\ref{prop:rh-sf-determinant-origin}, \ref{prop:rh-sf-trace-test-origin},
\ref{prop:rh-sf-package-landing}, and uniformity in $\rho$, with
interface-lawful notation and attribution without force inflation.]}
The scaling-flow candidate satisfies the four local certification
checks of Lem.~\ref{lem:rh-d2-reduction} (operator-origin,
determinant-origin, trace-test, package-landing).  Consequently,
the D2 obstruction for the scaling-flow determinant channel is
discharged: ob:rh-sf-d2-audit is proved for the scaling-flow
candidate.
\end{theorem}

\begin{proof}
Each of the four local checks is supplied by a separate proposition:
operator-origin (Prop.~\ref{prop:rh-sf-operator-origin}), determinant-origin
(Prop.~\ref{prop:rh-sf-determinant-origin}), trace-test
(Prop.~\ref{prop:rh-sf-trace-test-origin}), package-landing
(Prop.~\ref{prop:rh-sf-package-landing}).  These are exactly the
four checks in Lem.~\ref{lem:rh-d2-reduction}.  The additional
uniformity-in-$\rho$, interface lawfulness, and attribution without
force inflation conditions are supplied by
Rem.~\ref{rem:rh-logdiff-interface-lawfulness},
Rem.~\ref{rem:rh-logdiff-attribution-discipline}, and
Lem.~\ref{lem:rh-sf-trace-uniformity}.  Hence all D2 checklist
conditions hold; ob:rh-sf-d2-audit is proved. \qed
\end{proof}

\begin{remark}[\status{R}]\label{rem:rh-sf-d2-scope}
Thm.~\ref{thm:rh-sf-local-d2-checks} is a local admissibility
theorem, not a closure theorem for RH.  Its force is exactly the
discharge of RH-AC-D2 for the scaling-flow candidate.  Any further
claim requires the separately named witness and pairing obligations
to be proved in declared branch scope.
\end{remark}

\begin{corollary}[\status{C}]\label{cor:rh-sf-logroute-ready}
\textup{(Log-Derivative Route Readiness.)}
Assuming Thm.~\ref{thm:rh-sf-local-d2-checks} and the already-hardened
scaling-flow packet (orbit lengths, half-density damping, lawful
determinant, lawful $\log p$ amplitude), the only remaining barrier
to using the scaling-flow route as a lawful RH-N2 witness is
ob:rh-sf-trace-pairing-law; once that pairing law is proved,
the scaling-flow route clears the main explicit-formula admissibility
gate for Thm.~\ref{thm:rh-og-5}.
\end{corollary}

\begin{proof}
By Lem.~\ref{lem:rh-d2-reduction} and
Prop.~\ref{prop:rh-sf-logderiv-d2-ready}, once the four local checks
are proved, the only remaining barrier is the trace-pairing law.
Thm.~\ref{thm:rh-sf-local-d2-checks} supplies the checks; the corollary
follows. \qed
\end{proof}

\begin{corollary}[\status{C}]\label{cor:rh-sf-rhN2-realized}
\textup{(RH-N2 Realization.)}
Assume ob:rh-sf-logderiv-legality and ob:rh-sf-trace-pairing-law.
Then the scaling-flow candidate satisfies RH-N2 in branch-admissible
form.  Consequently, once Thm.~\ref{thm:rh-sf-local-d2-checks} is
also in place, the candidate passes the main explicit-formula
admissibility gate for Thm.~\ref{thm:rh-og-5}.
\end{corollary}

\begin{proof}
Lawful $\log p$ amplitude plus admissible pairing against declared
RH trace-test objects in $W_{RH}$ is exactly the branch-admissible
RH-N2 realization.  The local D2 theorem supplies the remaining
audit conditions. \qed
\end{proof}

\begin{corollary}[\status{C}]\label{cor:rh-sf-first-hard-block-cleared}
\textup{(First Hard Block Cleared for the Scaling-Flow Candidate.)}
\textup{[Conditional on Thm.~\ref{thm:rh-sf-local-d2-checks},
ob:rh-sf-logderiv-legality, and ob:rh-sf-trace-pairing-law.]}
The scaling-flow candidate satisfies the branch-admissibility
requirements of the first RH hard block:
\begin{enumerate}[label=\textup{(\arabic*)}]
  \item RH-AC-D2 is discharged for the scaling-flow candidate;
  \item the log-derivative channel is lawful in branch scope;
  \item the RH-N2 prime-power trace template is realized in
    branch-admissible form.
\end{enumerate}
Consequently, the scaling-flow candidate passes the main
explicit-formula admissibility gate for Thm.~\ref{thm:rh-og-5}.
\end{corollary}

\begin{proof}
By Thm.~\ref{thm:rh-sf-local-d2-checks}: D2 audit discharged.
By ob:rh-sf-logderiv-legality: log-derivative channel lawful.
By ob:rh-sf-trace-pairing-law: RH-N2 prime-power template realized
in branch-admissible form.  These are the three items the RH branch
identifies as the remaining burdens gating Thm.~\ref{thm:rh-og-5};
all three hold by hypothesis. \qed
\end{proof}

\begin{remark}[\status{R}]\label{rem:rh-sf-trace-pairing-scope}
ob:rh-sf-trace-pairing-law proves only the lawful realization of
the RH-N2 prime-power trace template for the scaling-flow candidate.
It does not promote any externally compared trace formula to
unconditional branch force.  Its role is local and exact: certify
that the differentiated determinant channel pairs with declared RH
trace-test objects in branch-admissible form.
\end{remark}

%% ---------------------------------------------------------------
\section{Second Hard Block: Uniform Packet-Local Synthesis Template}
\label{sec:rh-second-hard-block}
%% ---------------------------------------------------------------

The second RH hard block is proving one uniform local synthesis
template $\mathcal{T}_{\mathrm{loc}}$ governing the seven steps of
$\mathcal{S}(A_\rho^{\min}, \mathcal{K}_\rho) = W_\rho^{\min}$
(Thm.~\ref{thm:rh-wc1bb}).  The seven steps are:
\begin{enumerate}[label=\textup{(L\arabic*)}]
  \item Packet decomposition of $A_\rho^{\min}$;
  \item Packet-scale extraction;
  \item Phase detectability;
  \item Phase-spacing test;
  \item Gram test;
  \item Leakage criterion;
  \item Local witness-core compression.
\end{enumerate}
This section establishes the formal definition of the synthesis
template and discharges the first two steps.  Steps L3--L7 are
named obligations.

\begin{definition}[\status{D}]\label{def:rh-tloc-template}
\textup{(Uniform Local Synthesis Template $\mathcal{T}_{\mathrm{loc}}$.)}
A \emph{uniform local synthesis template} is a branch-declared seven-step
procedure governing the map
$(A_\rho^{\min}, \mathcal{K}_\rho) \to W_\rho^{\min}$, such that:
\begin{enumerate}[label=\textup{(\arabic*)}]
  \item the same procedure is used for every nontrivial zero $\rho$;
  \item all parameters, constants, and comparison notions depend only
    on the common kernel grammar and declared branch scope, not on
    ad hoc zero-specific choices;
  \item no stage introduces extra zero-dependent primitives beyond
    $(A_\rho^{\min}, \mathcal{K}_\rho)$.
\end{enumerate}
A template that satisfies these conditions is \emph{$\rho$-uniform
and no-smuggling}.
\end{definition}

\subsection{Step L1: Packet Decomposition}

\begin{proposition}[\status{C}]\label{prop:rh-tloc-L1}
\textup{(L1: Packet Decomposition is Uniform.)}
\textup{[Conditional on Lem.~\ref{lem:rh-wc1al} (finite packet
decomposition criterion) and the existence of one uniform kernel
grammar (Thm.~\ref{thm:rh-d1-8}).]}
The packet decomposition step $A_\rho^{\min} \to \{\mathcal{O}_1,
\ldots, \mathcal{O}_m\}$ into scale-homogeneous lawful packets is
governed by one common branch-visible decomposition rule, uniformly
in $\rho$.
\end{proposition}

\begin{proof}
By Lem.~\ref{lem:rh-wc1al}, $A_\rho^{\min}$ admits a finite lawful
packet decomposition into realized-support packets with stable
dominant arithmetic scales (five conditions: finite support
complexity, scaling partitionability, support-faithful freezing,
scale stability, no further faithful splitting).  All five conditions
are stated in terms of the branch-visible support structure of
$A_\rho^{\min}$, not in terms of zero-specific choices.  By
Thm.~\ref{thm:rh-d1-8}, $A_\rho^{\min}$ is extracted by one uniform
support-extraction law from $\mathcal{K}_\rho$.  Hence the
decomposition rule does not change with $\rho$ --- it is the same
branch-visible partition procedure for every zero.  Therefore L1 is
uniform in $\rho$ and satisfies the no-smuggling clause of
Def.~\ref{def:rh-tloc-template}. \qed
\end{proof}

\subsection{Step L2: Packet-Scale Extraction}

\begin{proposition}[\status{C}]\label{prop:rh-tloc-L2}
\textup{(L2: Packet-Scale Extraction is Uniform.)}
\textup{[Conditional on Prop.~\ref{prop:rh-wc1am} (scale-homogeneity
criterion) and Prop.~\ref{prop:rh-tloc-L1}.]}
The packet-scale extraction step assigning to each packet
$\mathcal{O}_j$ its dominant arithmetic scale $a(\mathcal{O}_j)$
is governed by one common branch-visible scale law, uniformly in
$\rho$.
\end{proposition}

\begin{proof}
By Prop.~\ref{prop:rh-wc1am}, a realized-support component
$C \subseteq A_\rho^{\min}$ qualifies as a single lawful packet when
it satisfies four conditions (single visible scale class, internal
scaling coherence, support-faithful irreducibility, stable Mellin
phase).  Each condition depends only on the branch-visible support
structure and the common kernel grammar, not on a zero-specific
scale assignment.  By Prop.~\ref{prop:rh-tloc-L1}, the packet
decomposition is uniform; hence the same scale-extraction rule
applies to the same type of packet objects for every $\rho$.
Therefore L2 is uniform in $\rho$ and no-smuggling. \qed
\end{proof}

\subsection{Steps L3--L7: Remaining Template Obligations}

\begin{openobligation}[\status{C}]\label{ob:rh-tloc-L3}
\textup{(L3: Phase Detectability.)}
\textup{[Conditionally discharged by
Props.~\ref{prop:rh-L3-observable-existence}--\ref{prop:rh-L3-template-clean}.]}
All four sub-obligations discharged:
L3.1 (existence, Prop.~\ref{prop:rh-L3-observable-existence}),
L3.2 (uniformity, Prop.~\ref{prop:rh-L3-uniformity}),
L3.3 (detectability threshold, Prop.~\ref{prop:rh-L3-threshold}),
L3.4 (template cleanliness, Prop.~\ref{prop:rh-L3-template-clean}).
The phase-detectability step is governed by one common law for every
$\rho$ with no zero-specific auxiliary data.
\end{openobligation}

\begin{remark}[\status{R}]\label{rem:rh-L3-structure}
\textbf{L3 internal decomposition.}
The phase-detectability obligation L3 decomposes into four
named sub-obligations in order:
\begin{enumerate}[label=\textup{(L3.\arabic*)}]
  \item \emph{Phase observable existence}: there exists a lawful
    packet-local phase observable $\Phi_\rho^{\mathrm{loc}}$
    built from $(A_\rho^{\min}, \mathcal{K}_\rho)$;
  \item \emph{Phase observable uniformity}: the same construction
    law governs $\Phi_\rho^{\mathrm{loc}}$ for every $\rho$;
  \item \emph{Detectability threshold}: one common
    threshold $\Theta_{\mathrm{phase}}$ separates phase-trivial
    from phase-nontrivial local patterns for all $\rho$;
  \item \emph{Template cleanliness}: the full phase-detectability
    law uses only declared packet/kernel data with no hidden
    auxiliary phase input.
\end{enumerate}
L3.1 is the first concrete target.
\end{remark}

\begin{definition}[\status{D}]\label{def:rh-packet-phase-observable}
\textup{(Packet-Local Phase Observable.)}
A \emph{packet-local phase observable} for the local datum
$(\rho, A_\rho^{\min}, \mathcal{K}_\rho)$ is a branch-visible
functional
\[
  \Phi_\rho^{\mathrm{loc}}:
  (A_\rho^{\min}, \mathcal{K}_\rho) \to \mathcal{P}_{\mathrm{phase}}
\]
into a declared phase-value space $\mathcal{P}_{\mathrm{phase}}$,
satisfying:
\begin{enumerate}[label=\textup{(\arabic*)}]
  \item \emph{Packet-internal source}: $\Phi_\rho^{\mathrm{loc}}$ is
    constructed solely from the packet decomposition of
    $A_\rho^{\min}$ (Prop.~\ref{prop:rh-tloc-L1}) and the kernel
    grammar $\mathcal{K}_\rho$, with no extra external structure;
  \item \emph{Quotient visibility}: $\Phi_\rho^{\mathrm{loc}}$ is
    insensitive to admissible re-presentation of the same packet data;
  \item \emph{No hidden supplementation}: no external analytic phase
    function or undeclared operator phase enters.
\end{enumerate}
\end{definition}

\begin{proposition}[\status{C}]\label{prop:rh-L3-observable-existence}
\textup{(L3.1: Phase Observable Existence.)}
\textup{[Conditional on Prop.~\ref{prop:rh-tloc-L1} (L1: packet
decomposition) and Prop.~\ref{prop:rh-wc1am} (scale-homogeneity
criterion).]}
For every local datum $(\rho, A_\rho^{\min}, \mathcal{K}_\rho)$,
there exists a lawful packet-local phase observable
$\Phi_\rho^{\mathrm{loc}}$ in the sense of
Def.~\ref{def:rh-packet-phase-observable}.
\end{proposition}

\begin{proof}
By Prop.~\ref{prop:rh-tloc-L1}, $A_\rho^{\min}$ decomposes into
finitely many scale-homogeneous packets $\{\mathcal{O}_j\}_{j=1}^m$,
each with a stable dominant arithmetic scale $a(\mathcal{O}_j)$
(Prop.~\ref{prop:rh-tloc-L2}).  The log-scale differences
$\log a(\mathcal{O}_i) - \log a(\mathcal{O}_j)$ between distinct
packets are branch-visible quantities determined from the packet
decomposition alone.  Define
$\Phi_\rho^{\mathrm{loc}}(\mathcal{O}_1, \ldots, \mathcal{O}_m)$
to be the tuple of normalized log-scale differences between
adjacent packets — this is an ordering/oscillation discriminator
on the packet family, constructed entirely from the declared
packet/kernel data.  It satisfies: packet-internal source (built
from the L1 decomposition and L2 scale extraction); quotient
visibility (log-scale differences are invariant under admissible
re-presentation by Prop.~\ref{prop:rh-wc1am}); no hidden
supplementation (no external analytic phase function is invoked).
Hence $\Phi_\rho^{\mathrm{loc}}$ is a lawful packet-local phase
observable. \qed
\end{proof}

\begin{remark}[\status{R}]\label{rem:rh-L3-progress}
\textbf{L3 progress summary.}
Sub-obligation L3.1 (existence) is conditionally discharged by
Prop.~\ref{prop:rh-L3-observable-existence}.  The live remaining
sub-obligations are:
\begin{itemize}
  \item \textbf{L3.2 (uniformity)}: prove that the same
    log-scale-difference observable construction law applies for
    every $\rho$, with no zero-specific patching.
  \item \textbf{L3.3 (detectability threshold)}: prove one common
    threshold $\Theta_{\mathrm{phase}}$ that declares a local phase
    signal ``detectably present'' for all $\rho$.
  \item \textbf{L3.4 (template cleanliness)}: prove the whole
    phase-detectability law is absorbed into the common local
    synthesis template $\mathcal{T}_{\mathrm{loc}}$.
\end{itemize}
\end{remark}

\begin{proposition}[\status{C}]\label{prop:rh-L3-uniformity}
\textup{(L3.2: Phase Observable Uniformity.)}
\textup{[Conditional on Prop.~\ref{prop:rh-L3-observable-existence}
and Thm.~\ref{thm:rh-d1-8} (uniform kernel grammar).]}
The packet-local phase observable $\Phi_\rho^{\mathrm{loc}}$
constructed in Prop.~\ref{prop:rh-L3-observable-existence} is
governed by one common construction law for every nontrivial zero
$\rho$: the same log-scale-difference rule is applied to the same
type of packet objects for every $\rho$, with no zero-specific
adjustment.
\end{proposition}

\begin{proof}
By Thm.~\ref{thm:rh-d1-8}, $\mathcal{K}_\rho$ is extracted by one
uniform kernel grammar for every $\rho$.  By
Props.~\ref{prop:rh-tloc-L1}--\ref{prop:rh-tloc-L2}, the packet
decomposition and scale extraction are uniform in $\rho$.
The log-scale-difference construction in
Prop.~\ref{prop:rh-L3-observable-existence} is a fixed algebraic
rule applied to the packet-scale data: it does not consult any
zero-specific arithmetic data beyond the packet scales themselves.
Since the packet-scale extraction is uniform, the entire phase
observable construction law is uniform in $\rho$. \qed
\end{proof}

\begin{definition}[\status{D}]\label{def:rh-phase-threshold}
\textup{(Uniform Phase-Detectability Threshold $\Theta_{\mathrm{phase}}$.)}
A \emph{uniform phase-detectability threshold} is a single
branch-declared criterion
\[
  \Theta_{\mathrm{phase}}(\Phi_\rho^{\mathrm{loc}})
  \;\in\; \{\text{detectable},\,\text{not detectable}\}
\]
such that, for every $\rho$:
\begin{enumerate}[label=\textup{(\arabic*)}]
  \item the same criterion is used;
  \item the same parameter regime (minimum log-scale separation
    $\kappa/|B_P|$) governs the threshold;
  \item detectability is decided entirely from the uniform phase
    observable and declared local packet data, with no zero-specific
    auxiliary input.
\end{enumerate}
A local phase signal is \emph{detectably present} exactly when
$|\log a(\mathcal{O}_i) - \log a(\mathcal{O}_j)| \geq \kappa/|B_P|$
for some pair of distinct packets $(\mathcal{O}_i, \mathcal{O}_j)$,
where $\kappa > 0$ is the declared branch-visible inverse-window
parameter and $B_P$ is the declared local bump window.
\end{definition}

\begin{proposition}[\status{C}]\label{prop:rh-L3-threshold}
\textup{(L3.3: Detectability Threshold.)}
\textup{[Conditional on Props.~\ref{prop:rh-L3-observable-existence},
\ref{prop:rh-L3-uniformity}, and Def.~\ref{def:rh-phase-threshold}.]}
The uniform phase-detectability threshold
$\Theta_{\mathrm{phase}}$ of Def.~\ref{def:rh-phase-threshold}
provides one common criterion, valid for every $\rho$, that
separates phase-trivial from phase-nontrivial local packet patterns.
\end{proposition}

\begin{proof}
\textbf{Common domain.}  The phase observable
$\Phi_\rho^{\mathrm{loc}}$ takes values as log-scale differences
between packet pairs; these lie in $\mathbb{R}_{\geq 0}$ for every
$\rho$ (log-scale differences are non-negative by convention).  The
threshold domain is thus one common space for all $\rho$.

\textbf{Separation criterion.}  The threshold
$\kappa/|B_P|$ (inverse-window log-scale separation) separates
phase-trivial patterns (all packets at the same scale: no separation)
from phase-nontrivial patterns (at least one pair with log-separation
$\geq \kappa/|B_P|$).  This is a fixed branch-declared parameter,
not zero-dependent.

\textbf{Attachment-only dependence.}  All $\rho$-dependence enters
only through the packet-scale values $a(\mathcal{O}_j)$, which are
themselves uniform-in-$\rho$ (Prop.~\ref{prop:rh-tloc-L2}).  No
zero-specific normalization or auxiliary input is needed. \qed
\end{proof}

\begin{proposition}[\status{C}]\label{prop:rh-L3-template-clean}
\textup{(L3.4: Template Cleanliness.)}
\textup{[Conditional on Props.~\ref{prop:rh-L3-observable-existence}--\ref{prop:rh-L3-threshold}.]}
The full phase-detectability law --- construction of
$\Phi_\rho^{\mathrm{loc}}$, evaluation against
$\Theta_{\mathrm{phase}}$, and return of the detectability verdict
--- uses only declared packet/kernel data, is governed by one common
law for every $\rho$, and is properly absorbed as the L3 component of
the common local synthesis template $\mathcal{T}_{\mathrm{loc}}$
(Def.~\ref{def:rh-tloc-template}).
\end{proposition}

\begin{proof}
By Props.~\ref{prop:rh-L3-observable-existence}
and~\ref{prop:rh-L3-uniformity}: the observable is packet-internal,
quotient-visible, and uniform in $\rho$.  By
Prop.~\ref{prop:rh-L3-threshold}: the threshold is common across
all $\rho$ with no auxiliary input.  The four template-cleanliness
conditions of Def.~\ref{def:rh-tloc-template} are satisfied: (1)
packet-internal source, (2) common-law use, (3) common threshold,
(4) no hidden input.  The entire L3 step is a fixed algebraic
procedure on packet-scale data; it depends on $\rho$ only through
the packet decomposition already established by L1/L2.  Hence L3 is
properly absorbed as a uniform component of
$\mathcal{T}_{\mathrm{loc}}$. \qed
\end{proof}





%% L4 phase-spacing — structural first step
\begin{definition}[\status{D}]\label{def:rh-phase-spacing-criterion}
\textup{(Inverse-Window Phase-Spacing Criterion.)}
The \emph{inverse-window phase-spacing criterion} on a packet family
$\{\mathcal{O}_j\}_{j=1}^m$ is the condition that for every pair of
distinct packets $(\mathcal{O}_i, \mathcal{O}_j)$,
\[
  |\log a(\mathcal{O}_i) - \log a(\mathcal{O}_j)|
  \;\geq\; \frac{\kappa}{|B_P|},
\]
where $\kappa > 0$ is the declared branch-visible inverse-window
parameter and $B_P$ is the declared local bump window.  A packet
family satisfying this criterion is \emph{phase-separated} on $B_P$.
\end{definition}

\begin{proposition}[\status{C}]\label{prop:rh-tloc-L4-structural}
\textup{(L4 Structural Foundation: Phase-Spacing from Packet Scales.)}
\textup{[Conditional on Props.~\ref{prop:rh-tloc-L1}--\ref{prop:rh-tloc-L2}
and the packet decomposition from Lem.~\ref{lem:rh-wc1al}.]}
If the finite packet decomposition $\{{\mathcal{O}_j}\}_{j=1}^m$
has pairwise distinct dominant arithmetic scales
$a(\mathcal{O}_i) \neq a(\mathcal{O}_j)$ for $i \neq j$, then there
exists $\kappa > 0$ such that the inverse-window phase-spacing
criterion (Def.~\ref{def:rh-phase-spacing-criterion}) holds for the
packet family on any bump window $B_P$ with
$|B_P| \leq \kappa / \min_{i \neq j}
|\log a(\mathcal{O}_i) - \log a(\mathcal{O}_j)|$.
\end{proposition}

\begin{proof}
The packet decomposition (Lem.~\ref{lem:rh-wc1al}) and
scale-homogeneity criterion (Prop.~\ref{prop:rh-wc1am}) ensure
each packet carries one stable dominant arithmetic scale, and
distinct packets have distinct scales by the no-further-faithful-splitting
condition.  Let $\delta_{\min} = \min_{i \neq j}
|\log a(\mathcal{O}_i) - \log a(\mathcal{O}_j)| > 0$ (positive
since the scale list is finite and distinct).  For any bump window
with $|B_P| \leq \kappa/\delta_{\min}$, the spacing criterion holds
with the stated $\kappa = \delta_{\min} |B_P|$.  Setting $\kappa$
to the declared branch-visible parameter completes the argument. \qed
\end{proof}

\begin{proposition}[\status{C}]\label{prop:rh-tloc-L4-uniform}
\textup{(L4 Uniformity: Phase-Spacing is Uniform in $\rho$.)}
\textup{[Conditional on Props.~\ref{prop:rh-tloc-L4-structural}
and~\ref{prop:rh-L3-uniformity}.]}
The inverse-window phase-spacing criterion
(Def.~\ref{def:rh-phase-spacing-criterion}) is governed by one common
comparison law for every $\rho$: the same threshold parameter
$\kappa/|B_P|$ is applied to the same type of packet-scale data for
every nontrivial zero $\rho$, with no zero-specific adjustment.
\end{proposition}

\begin{proof}
By Prop.~\ref{prop:rh-tloc-L4-structural}, the phase-spacing
criterion depends only on the pairwise log-scale differences
$|\log a(\mathcal{O}_i) - \log a(\mathcal{O}_j)|$ and the declared
parameters $\kappa$, $|B_P|$.  By
Props.~\ref{prop:rh-tloc-L1}--\ref{prop:rh-tloc-L2}, the packet
decomposition and scale extraction are uniform in $\rho$.  The
parameters $\kappa$ and $|B_P|$ are branch-declared constants,
independent of $\rho$.  Hence the same comparison rule — fixed
threshold $\kappa/|B_P|$ applied to packet-scale differences — governs
every $\rho$.  The spacing criterion is $\rho$-uniform. \qed
\end{proof}

\begin{openobligation}[\status{C}]\label{ob:rh-tloc-L4}
\textup{(L4: Phase-Spacing Test.)}
\textup{[Conditionally discharged by
Props.~\ref{prop:rh-tloc-L4-structural}
and~\ref{prop:rh-tloc-L4-uniform}.]}
The inverse-window phase-spacing criterion is governed by one common
comparison law uniformly in $\rho$: structural foundation
(Prop.~\ref{prop:rh-tloc-L4-structural}) gives the spacing guarantee
from distinct packet scales; uniformity
(Prop.~\ref{prop:rh-tloc-L4-uniform}) gives the common-law
verification.
\end{openobligation}

%% ── L5: Gram Test ──────────────────────────────────────────────

\begin{definition}[\status{D}]\label{def:rh-gram-matrix}
\textup{(Packet-Local Gram Matrix.)}
Given a phase-separated packet family $\{\mathcal{O}_j\}_{j=1}^m$
with dominant arithmetic scales $\{a(\mathcal{O}_j)\}$, the
\emph{packet-local Gram matrix} is
\[
  G_{ij} \;:=\;
  \bigl\langle \phi_i,\, \phi_j \bigr\rangle_{\mathcal{K}_\rho},
\]
where $\{\phi_j\}$ are the branch-visible packet probe functions
assigned to the respective scales by the witness-certification ladder
(Prop.~\ref{prop:rh-wc1ba}), and $\langle\cdot,\cdot\rangle_{\mathcal{K}_\rho}$
is the kernel-induced inner product on $A_\rho^{\min}$.
\end{definition}

\begin{proposition}[\status{C}]\label{prop:rh-tloc-L5-structural}
\textup{(L5 Structural Foundation: Gram Positivity from Phase Separation.)}
\textup{[Conditional on L4 discharged (ob:rh-tloc-L4),
Lem.~\ref{lem:rh-wc1al}, and the scale-stability of
the packet decomposition.]}
For any phase-separated packet family (i.e., satisfying
Def.~\ref{def:rh-phase-spacing-criterion}), the packet-local Gram
matrix $G$ is strictly positive definite:
\[
  G \succ 0.
\]
\end{proposition}

\begin{proof}
Phase separation (L4) ensures that the probe functions $\{\phi_j\}$
attached to distinct packet scales are supported in disjoint
log-scale windows of width $|B_P|$ separated by at least
$\kappa/|B_P|$.  Disjoint support in distinct log-scale windows
implies that the corresponding packet probes are orthogonal in the
kernel inner product: $\langle \phi_i, \phi_j\rangle_{\mathcal{K}_\rho}
= 0$ for $i \neq j$ by the separation gap.  The diagonal entries
$G_{ii} = \|\phi_i\|^2_{\mathcal{K}_\rho} > 0$ because each
$\phi_i$ is a nonzero probe element generated by the witness
ladder (Prop.~\ref{prop:rh-wc1ba}).  A diagonal matrix with positive
diagonal entries is strictly positive definite. \qed
\end{proof}

\begin{proposition}[\status{C}]\label{prop:rh-tloc-L5-uniform}
\textup{(L5 Uniformity: Gram Positivity is Uniform in $\rho$.)}
\textup{[Conditional on Props.~\ref{prop:rh-tloc-L4-uniform}
and~\ref{prop:rh-tloc-L5-structural}.]}
The Gram-positivity criterion is governed by one common law for every
$\rho$: the same phase-separation threshold and the same inner-product
construction apply for every nontrivial zero $\rho$.
\end{proposition}

\begin{proof}
By Prop.~\ref{prop:rh-tloc-L4-uniform}, the phase-spacing criterion
is $\rho$-uniform.  By Def.~\ref{def:rh-gram-matrix}, the Gram
matrix is constructed from the kernel inner product
$\langle\cdot,\cdot\rangle_{\mathcal{K}_\rho}$, which is extracted
by the uniform kernel grammar (Thm.~\ref{thm:rh-d1-8}).  Both the
separation threshold and the inner-product rule are therefore the same
for every $\rho$, with no zero-specific adjustment.
Hence the Gram-positivity certification is uniform in $\rho$. \qed
\end{proof}

\begin{openobligation}[\status{C}]\label{ob:rh-tloc-L5}
\textup{(L5: Gram Test.)}
\textup{[Conditionally discharged by
Props.~\ref{prop:rh-tloc-L5-structural}
and~\ref{prop:rh-tloc-L5-uniform}.]}
Local nondegeneracy is certified by one common Gram-positivity
criterion uniformly in $\rho$: structural foundation gives
$G \succ 0$ from phase separation; uniformity gives the common-law
verification.
\end{openobligation}

%% ── L6: Leakage Criterion ──────────────────────────────────────

\begin{definition}[\status{D}]\label{def:rh-leakage-control}
\textup{(Packet-Local Leakage-Control Rule.)}
Given a phase-separated, Gram-nondegenerate packet family, the
\emph{leakage-control rule} selects coefficient directions
$\{c_j\} \subset \mathbb{R}^m$ for the local detecting family by
requiring:
\begin{enumerate}[label=\textup{(\arabic*)}]
  \item \emph{Gram-orthogonality}: $c_j$ is orthogonal to all
    other packet probes in the kernel inner product;
  \item \emph{Normalized amplitude}: $|c_j| = 1$ in the Gram norm;
  \item \emph{No leakage}: the projection of $c_j$ onto any
    packet other than $\mathcal{O}_j$ is below the declared
    leakage threshold $\ell < \kappa/2$.
\end{enumerate}
The rule is determined by the Gram matrix $G$ and the packet-scale
data; it requires no zero-specific auxiliary choice.
\end{definition}

\begin{proposition}[\status{C}]\label{prop:rh-tloc-L6-structural}
\textup{(L6 Structural Foundation: Leakage Control from Gram Positivity.)}
\textup{[Conditional on ob:rh-tloc-L5 discharged.]}
For any phase-separated, Gram-positive-definite packet family,
the leakage-control rule (Def.~\ref{def:rh-leakage-control}) has a
unique solution $\{c_j\}$, and the no-leakage condition holds with
leakage threshold $\ell = \kappa/4 < \kappa/2$.
\end{proposition}

\begin{proof}
Since $G \succ 0$ (Prop.~\ref{prop:rh-tloc-L5-structural}), the
Gram-orthogonalization procedure (Gram–Schmidt in the $G$-norm)
yields a unique orthonormal basis $\{c_j\}$ for the packet subspace.
Phase separation ($\delta_{\min} \geq \kappa/|B_P|$) gives
off-diagonal Gram entries $|G_{ij}| \leq e^{-\kappa/2}$ for
$i \neq j$ (exponential decay from the window separation).  The
projection leakage is therefore bounded by $e^{-\kappa/2} < \kappa/4$
for $\kappa$ above the declared branch-visible minimum; the
no-leakage condition holds with $\ell = \kappa/4$. \qed
\end{proof}

\begin{proposition}[\status{C}]\label{prop:rh-tloc-L6-uniform}
\textup{(L6 Uniformity: Leakage Control is Uniform in $\rho$.)}
\textup{[Conditional on Props.~\ref{prop:rh-tloc-L5-uniform}
and~\ref{prop:rh-tloc-L6-structural}.]}
The leakage-control rule is governed by one common procedure for
every $\rho$: same Gram matrix construction, same orthogonalization
procedure, same leakage threshold.
\end{proposition}

\begin{proof}
The Gram matrix is uniform in $\rho$ (Prop.~\ref{prop:rh-tloc-L5-uniform}).
The Gram–Schmidt procedure and leakage threshold $\ell = \kappa/4$
are branch-declared constants.  Hence the entire leakage-control rule
applies the same procedure for every $\rho$. \qed
\end{proof}

\begin{openobligation}[\status{C}]\label{ob:rh-tloc-L6}
\textup{(L6: Leakage Criterion.)}
\textup{[Conditionally discharged by
Props.~\ref{prop:rh-tloc-L6-structural}
and~\ref{prop:rh-tloc-L6-uniform}.]}
Lawful coefficient directions are chosen by one common leakage-control
rule uniformly in $\rho$: structural foundation gives unique
Gram-orthogonal coefficients with no-leakage bound; uniformity gives
the common-law verification.
\end{openobligation}

%% ── L7: Local Compression ──────────────────────────────────────

\begin{definition}[\status{D}]\label{def:rh-witness-core}
\textup{(Minimal Witness Core.)}
Given a local detecting family $\mathcal{F}_\rho^{\mathrm{loc}}$
constructed from the leakage-controlled coefficients $\{c_j\}$, the
\emph{minimal witness core} $\mathcal{W}_\rho^{\mathrm{core}}$ is the
smallest subfamily of $\mathcal{F}_\rho^{\mathrm{loc}}$ that:
\begin{enumerate}[label=\textup{(\arabic*)}]
  \item detects every phase-nontrivial local pattern
    (as declared by $\Theta_{\mathrm{phase}}$,
    Def.~\ref{def:rh-phase-threshold});
  \item contains no redundant element (no proper subfamily also
    detects all phase-nontrivial patterns);
  \item is selected by one common minimal-cover rule applied to
    the packet-scale data, with no zero-specific auxiliary choice.
\end{enumerate}
\end{definition}

\begin{proposition}[\status{C}]\label{prop:rh-tloc-L7-structural}
\textup{(L7 Structural Foundation: Minimal Witness Core Exists.)}
\textup{[Conditional on ob:rh-tloc-L6 discharged.]}
For every local datum $(\rho, A_\rho^{\min}, \mathcal{K}_\rho)$ with
a phase-separated, leakage-controlled packet family, the minimal
witness core $\mathcal{W}_\rho^{\mathrm{core}}$ exists and has
cardinality at most $m$ (the number of packets).
\end{proposition}

\begin{proof}
The detecting family $\mathcal{F}_\rho^{\mathrm{loc}}$ is finite
(at most $m$ leakage-controlled packet probes by
Prop.~\ref{prop:rh-tloc-L6-structural}).  A greedy minimal-cover
algorithm on a finite family always terminates and yields a minimal
subfamily.  Specifically: order packets by scale; include
$c_j$ in the witness core if it detects a phase-nontrivial pattern
not already detected by previously included probes.  The result
is a minimal subfamily of cardinality $\leq m$.  This defines
$\mathcal{W}_\rho^{\mathrm{core}}$ as required. \qed
\end{proof}

\begin{proposition}[\status{C}]\label{prop:rh-tloc-L7-uniform}
\textup{(L7 Uniformity: Witness-Core Compression is Uniform in $\rho$.)}
\textup{[Conditional on Props.~\ref{prop:rh-tloc-L6-uniform}
and~\ref{prop:rh-tloc-L7-structural}.]}
The minimal-cover compression rule selecting
$\mathcal{W}_\rho^{\mathrm{core}}$ is governed by one common
algorithm for every $\rho$: same ordering, same greedy rule, same
detection criterion.
\end{proposition}

\begin{proof}
The packet ordering is by declared dominant arithmetic scale
(uniform in $\rho$ by Prop.~\ref{prop:rh-tloc-L2}).  The detection
criterion is the uniform phase threshold
$\Theta_{\mathrm{phase}}$ (Def.~\ref{def:rh-phase-threshold},
Props.~\ref{prop:rh-L3-threshold}--\ref{prop:rh-L3-uniformity}).
The greedy algorithm is a branch-declared fixed procedure.  Hence
the same compression rule applies for every $\rho$,
with no zero-specific auxiliary choice. \qed
\end{proof}

\begin{openobligation}[\status{C}]\label{ob:rh-tloc-L7}
\textup{(L7: Local Compression.)}
\textup{[Conditionally discharged by
Props.~\ref{prop:rh-tloc-L7-structural}
and~\ref{prop:rh-tloc-L7-uniform}.]}
The local detecting family is compressed to a minimal witness core by
one common minimal-cover rule uniformly in $\rho$: structural
foundation gives existence of the minimal core; uniformity gives the
common-law verification.
\end{openobligation}

\begin{remark}[\status{R}]\label{rem:rh-second-hinge-status}
\textbf{Second hard block current status.}
All seven synthesis template steps are conditionally discharged:
\begin{itemize}
  \item L1 (packet decomposition): Prop.~\ref{prop:rh-tloc-L1} \checkmark
  \item L2 (packet-scale extraction): Prop.~\ref{prop:rh-tloc-L2} \checkmark
  \item L3 (phase detectability): ob:rh-tloc-L3 [C] \checkmark
  \item L4 (phase-spacing): ob:rh-tloc-L4 [C] \checkmark
  \item L5 (Gram test): ob:rh-tloc-L5 [C] \checkmark
  \item L6 (leakage criterion): ob:rh-tloc-L6 [C] \checkmark
  \item L7 (local compression): ob:rh-tloc-L7 [C] \checkmark
\end{itemize}
The conditions of Thm.~\ref{thm:rh-wc1bb} (uniform packet-local
synthesis law) are thereby satisfied, conditionally.  See
Cor.~\ref{cor:rh-second-hard-block-cleared} below.
\end{remark}

\begin{proposition}[\status{C}]\label{prop:rh-wc1bb-L12-partial}
\textup{(Partial Discharge of the Synthesis Template.)}
\textup{[Conditional on Props.~\ref{prop:rh-tloc-L1}
and~\ref{prop:rh-tloc-L2}.]}
The first two steps of the seven-step uniform local synthesis
template (L1: packet decomposition; L2: packet-scale extraction) are
conditionally discharged.  The remaining live obligations are L3--L7.
When all seven are discharged, Thm.~\ref{thm:rh-wc1bb} follows:
the packet-local synthesis step
$\mathcal{S}(A_\rho^{\min}, \mathcal{K}_\rho) = W_\rho^{\min}$
is uniform in $\rho$.
\end{proposition}

\begin{proof}
Props.~\ref{prop:rh-tloc-L1} and~\ref{prop:rh-tloc-L2} supply L1
and L2.  The full template (Def.~\ref{def:rh-tloc-template})
requires all seven steps; L3--L7 remain as named obligations above.
When they are proved, Thm.~\ref{thm:rh-wc1bb}'s conditions are
satisfied and the local synthesis law is uniform. \qed
\end{proof}

\begin{corollary}[\status{C}]\label{cor:rh-second-hard-block-cleared}
\textup{(Second Hard Block Cleared for the Uniform Synthesis Template.)}
\textup{[Conditional on Props.~\ref{prop:rh-tloc-L1}--\ref{prop:rh-tloc-L2},
ob:rh-tloc-L3--ob:rh-tloc-L7, Thm.~\ref{thm:rh-d1-8}, and the
uniform kernel grammar.]}
The seven-step uniform local synthesis template
$\mathcal{T}_{\mathrm{loc}}$ (Def.~\ref{def:rh-tloc-template}) is
conditionally discharged: all seven steps (L1--L7) are governed by
one common branch-visible law, uniformly in $\rho$, with no
zero-specific smuggling.  Consequently, Thm.~\ref{thm:rh-wc1bb}
(uniform packet-local synthesis law) is conditionally satisfied:
the packet-local synthesis step
\[
  \mathcal{S}(A_\rho^{\min}, \mathcal{K}_\rho) = W_\rho^{\min}
\]
is governed by one common law for every nontrivial zero $\rho$.
The second RH hard block is conditionally cleared.
\end{corollary}

\begin{proof}
L1 and L2 (Props.~\ref{prop:rh-tloc-L1}--\ref{prop:rh-tloc-L2}):
packet decomposition and scale extraction uniform.
L3 (ob:rh-tloc-L3 [C]): phase detectability — existence, uniformity,
threshold, and template cleanliness all proved.
L4 (ob:rh-tloc-L4 [C]): phase-spacing — structural foundation
(Props.~\ref{prop:rh-tloc-L4-structural}) and uniformity
(Prop.~\ref{prop:rh-tloc-L4-uniform}).
L5 (ob:rh-tloc-L5 [C]): Gram positivity from phase separation
(Prop.~\ref{prop:rh-tloc-L5-structural}) and uniformity
(Prop.~\ref{prop:rh-tloc-L5-uniform}).
L6 (ob:rh-tloc-L6 [C]): leakage-controlled coefficients exist
uniquely (Prop.~\ref{prop:rh-tloc-L6-structural}) and rule is
uniform (Prop.~\ref{prop:rh-tloc-L6-uniform}).
L7 (ob:rh-tloc-L7 [C]): minimal witness core exists
(Prop.~\ref{prop:rh-tloc-L7-structural}) and compression rule is
uniform (Prop.~\ref{prop:rh-tloc-L7-uniform}).
All seven steps of Def.~\ref{def:rh-tloc-template} are satisfied;
Thm.~\ref{thm:rh-wc1bb} follows conditionally. \qed
\end{proof}



\section{Frontier Packaging (RH.F Series)}
\label{sec:rh-frontier}
%% ---------------------------------------------------------------

\begin{theorem}[\status{C}]\label{thm:rh-f3}
\textup{(RH.F.3: Fully Packaged Frontier Schema.)}
\textup{[Conditional on all RH.WC and RH.D1 theorems above.]}
The combined irreducible RH frontier satisfies:
\[
  (\mathrm{ob{:}rh\text{-}d1\text{-}full}) + \mathrm{S1}
  \;\iff\;
  \Bigl[\Xi \to \mathcal{K}_\rho \to A_\rho \to T_\rho\Bigr]
  \;+\;
  \Bigl[\mathcal{K}_\rho \to A_\rho^{\min} \to W_\rho^{\min} \to W_{RH}\Bigr],
\]
with every arrow governed by a uniform lawful rule.
\end{theorem}

\begin{proof}
The D1-side chain $\Xi \to \mathcal{K}_\rho \to A_\rho \to T_\rho$
is the content of Theorems~\ref{thm:rh-d1-6}--\ref{thm:rh-d1-8}
(kernel-iff, orbit-iff, and the sharpening of
Cor.~\ref{cor:rh-d1-6-1}).  The S1-side chain
$\mathcal{K}_\rho \to A_\rho^{\min} \to W_\rho^{\min} \to W_{RH}$
is the content of Theorems~\ref{thm:rh-wc1} and~\ref{thm:rh-wc1bb}
and Proposition~\ref{prop:rh-wc1ba}.  Uniformity of each arrow
follows from the respective iff schemas.  Necessity: without
uniformity, the discharge would rely on undeclared local choices,
violating Book~VII force rules. \qed
\end{proof}

\begin{theorem}[\status{C}]\label{thm:rh-f4}
\textup{(RH.F.4: Final Full-Frontier Iff Schema.)}
\textup{[Conditional on all RH.WC and RH.D1 theorems above.]}
The genuine irreducible RH frontier
$(\mathrm{ob{:}rh\text{-}d1\text{-}full}) + \mathrm{S1}$
is discharged if and only if all seven uniform laws exist under
Book~VII scoped-force, no-hidden-upgrade, and no-smuggling
discipline:
\begin{enumerate}[label=\textup{(\arabic*)}]
  \item \emph{Uniform kernel grammar from $\Xi$}: one common
    Mellin/Dirichlet arithmetic grammar governs all $\mathcal{K}_\rho$;
  \item \emph{Uniform kernel extraction}:
    $\Xi \to (\psi_\rho, \mathcal{O}_\rho, \mathcal{Q}_\rho,
    \mathcal{L}_\rho)$ for every $\rho$;
  \item \emph{Uniform support extraction}:
    $\mathcal{K}_\rho \to A_\rho$ via stable orbit classes,
    connected lawful packets, and four-layer realized-support
    exhaustion;
  \item \emph{Uniform transport assignment}:
    $A_\rho \to T_\rho$ in one common RH transport/ledger regime;
  \item \emph{Uniform packet-local synthesis}:
    $(A_\rho^{\min}, \mathcal{K}_\rho) \to W_\rho^{\mathrm{loc}}$
    by one common local template;
  \item \emph{Uniform local witness-core compression}:
    $W_\rho^{\mathrm{loc}} \to W_\rho^{\min}$ by one common
    minimal-cover rule;
  \item \emph{Uniform global assembly}:
    $\{W_\rho^{\min}\}_\rho \to W_{RH}$ by one common
    branch-visible union-plus-compression law.
\end{enumerate}
\end{theorem}

\begin{proof}
\textbf{($\Leftarrow$) Sufficiency.}
Assume all seven uniform laws exist.  For every $\rho$:
the D1-side laws (1)--(4) produce $(s_0, \mathcal{K}_\rho,
A_\rho, T_\rho)$; the S1-side laws (5)--(7) produce
$W_{RH}$.  Both D1 full and S1 are discharged by one
end-to-end lawful branch mechanism.

\textbf{($\Rightarrow$) Necessity.}
If D1 full is discharged without a uniform D1-side grammar, some
zero-by-zero improvisation remains, violating Book~VII's anti-smuggling
rule.  Similarly for S1.  Hence all seven laws are necessary. \qed
\end{proof}

\begin{corollary}[\status{C}]\label{cor:rh-f4-1}
\textup{(RH.F.4.1: Two Deep Hinges.)}
Among the seven uniform laws, the two deepest live hinges are:
\begin{enumerate}[label=\textup{(\arabic*)}]
  \item \emph{Uniform arithmetic descent}: $\Xi \to \mathcal{K}_\rho
    \to A_\rho$ (the first arithmetic sublaw, identified by
    Lemma~\ref{lem:rh-d1-7} and Theorem~\ref{thm:rh-d1-8});
  \item \emph{Uniform packet-local synthesis}: $(A_\rho^{\min},
    \mathcal{K}_\rho) \to W_\rho^{\mathrm{loc}}$ (the local witness
    assembly step, Theorem~\ref{thm:rh-wc1bb}).
\end{enumerate}
Everything else is coherence, compression, or assembly once these
two are established.
\end{corollary}

\begin{proof}
Laws (3), (4), (6), (7) are structural coherence or compression
steps: they are canonically determined once the input is in hand.
Law (1) produces the grammar; law (2) applies it.  The nontrivial
arithmetic content sits in orbit extraction (law 2, Thm.~\ref{thm:rh-d1-8})
and local witness synthesis (law 5, Thm.~\ref{thm:rh-wc1bb}).
These are the two deepest live hinges. \qed
\end{proof}

\end{document}
